% !TEX program = pdflatex
% !BIB program = bibtex
% Elsevier-compatible manuscript (IEEE twin: main.ltx). Computer Networks encourages the
% elsarticle class or the els-cas-templates bundle; see https://www.elsevier.com/latex
% Compile with Elsevier-provided elsarticle.cls.
% Two-column final layout: \documentclass[final,5p,times,twocolumn]{elsarticle},
% then (i) \renewenvironment{cps@table}{\begin{table}[!t]}{\end{table}}, (ii) restore figure*/table*
% for full-width floats. Preprint uses cps@table (pinned [H]; see preamble).
%
% References: this file uses BibTeX + elsarticle-num.bst (from the elsarticle bundle).
% If in-text cites show as [?], run the full chain from this directory:
%   pdflatex main_elsevier && bibtex main_elsevier && pdflatex main_elsevier && pdflatex main_elsevier
% Or: latexmk -pdf main_elsevier.tex (see latexmkrc in this folder).
%
% Pass natbib options before elsarticle loads natbib (avoids clashes with \biboptions alone).
\PassOptionsToPackage{numbers,sort&compress}{natbib}
\documentclass[preprint,12pt,times]{elsarticle}
% elsarticle: pass natbib options here too (avoids [?] if class did not inherit preamble-only \PassOptionsToPackage).
\biboptions{numbers,sort&compress}

% Uncomment for review PDFs with line numbers: \usepackage{lineno} and after \begin{document}: \linenumbers
\usepackage{amsmath,amssymb,amsfonts,amsthm}
\newtheorem{proposition}{Proposition}
\usepackage{algorithm}
\usepackage{algorithmic}
\usepackage{etoolbox}

% Tighter float glue ([H] tables/figures/algorithms should not add a full blank line at page top)
\setlength{\intextsep}{5pt plus 1pt minus 1pt}
\setlength{\textfloatsep}{6pt plus 1pt minus 1pt}
\setlength{\floatsep}{5pt plus 1pt minus 1pt}
\setlength{\abovecaptionskip}{4pt}
\setlength{\belowcaptionskip}{4pt}
\setlength{\headsep}{12pt}

% Algorithms: avoid horizontal spill and excess vertical list padding inside the ruled box
\AtBeginEnvironment{algorithm}{%
  \setlength{\topsep}{2pt}%
  \setlength{\partopsep}{0pt}%
}
\AtBeginEnvironment{algorithmic}{%
  \sloppy
  \setlength{\itemsep}{0pt}%
  \setlength{\parsep}{0pt}%
  \setlength{\topsep}{0pt}%
  \setlength{\partopsep}{0pt}%
}

\usepackage{graphicx}
\usepackage{adjustbox}
\usepackage{tikz}
\usetikzlibrary{shapes,arrows.meta,positioning,fit,calc,shadows,backgrounds}
\usepackage[section]{placeins} % keep floats within their section
\usepackage{float} % provides [H] to pin tables/figures in place (Overleaf-friendly)
% Do NOT load subcaption with elsarticle (caption compatibility mode → fatal on Overleaf).
% Multi-panel figures use minipage + (a)/(b) labels instead.
% Float placement: [H] keeps review PDFs predictable. For the final
% two-column layout (\documentclass[final,5p,times,twocolumn]{elsarticle}),
% \renewenvironment{cps@table}{\begin{table}[!t]}{\end{table}} and restore figure*/table*.
\newenvironment{cps@table}{\begin{table}[H]}{\end{table}}

% Allow LaTeX to mix text and floats on the same page (avoid one-figure float pages)
\renewcommand{\topfraction}{0.85}
\renewcommand{\bottomfraction}{0.75}
\renewcommand{\textfraction}{0.15}
\renewcommand{\floatpagefraction}{0.70}
\setcounter{topnumber}{3}
\setcounter{bottomnumber}{2}
\setcounter{totalnumber}{5}
% Reduce extra space at top of float-only pages
\makeatletter
\setlength{\@fptop}{0pt}
\setlength{\@fpbot}{0pt plus 1fil}
\makeatother

% Evaluation figures: scale by line width so on-page text matches Fig. 9/10.
% Do NOT cap height for most figures — height caps shrink axis labels in the PDF.
\newcommand{\evalfig}[1]{%
  \includegraphics[width=\linewidth]{#1}%
}
\newcommand{\evalfiglarge}[1]{%
  \includegraphics[width=\linewidth]{#1}%
}
\newcommand{\evalfigwide}[1]{%
  \includegraphics[width=\linewidth]{#1}%
}
\newcommand{\evalfigwideLarge}[1]{%
  \includegraphics[width=\linewidth]{#1}%
}
\newcommand{\evalfigpanel}[1]{%
  \includegraphics[width=\linewidth]{#1}%
}
\newcommand{\evalfigtall}[1]{%
  \includegraphics[width=\linewidth,height=0.52\textheight,keepaspectratio]{#1}%
}
\newcommand{\evalfigsmall}[1]{%
  \includegraphics[width=0.92\linewidth]{#1}%
}

\usepackage{tabularx} % full-width tables without resizebox
\usepackage{textcomp}
\usepackage{xcolor}
\usepackage{booktabs}
\usepackage{array}
\usepackage{url}
\usepackage{multirow}
\usepackage{makecell}
\usepackage[T1]{fontenc}
\usepackage[english]{babel}
\hyphenation{Zarbi Farzinvash Salehpour}
\usepackage{microtype} % improved justification and hyphenation
\usepackage[hidelinks]{hyperref}
% Uniform table typography for all floats:
% body = \footnotesize; notes = \scriptsize; never \resizebox (keeps font size consistent).
\newcommand{\tablebodyfont}{\footnotesize}
% Two-line CI cell: keeps numeric columns narrow so all tables share one font size.
% Keep the complete interval in one math expression so the comma renders correctly.
\newcommand{\cici}[3]{\makecell[c]{#1\\[-1pt]\scriptsize\ensuremath{[#2,\,#3]}}}
\newcommand{\cicib}[3]{\makecell[c]{\textbf{#1}\\[-1pt]\scriptsize\ensuremath{[#2,\,#3]}}}
\newcommand{\figref}[1]{\textbf{Fig.~\ref{#1}}}
\newcommand{\tabref}[1]{\textbf{Table~\ref{#1}}}
\newcommand{\evaltable}[1]{%
  \begingroup\tablebodyfont
  \renewcommand{\arraystretch}{1.15}%
  \setlength{\tabcolsep}{3.5pt}%
  #1%
  \endgroup
}
\newcommand{\tablenote}[1]{%
  \par\vspace{3pt}%
  {\scriptsize\raggedright #1\par}%
}
\newcommand{\tablenotecols}[2]{%
  \par\vspace{3pt}%
  {\scriptsize
  \begin{tabularx}{\linewidth}{@{}>{\raggedright\arraybackslash}X >{\raggedright\arraybackslash}X@{}}
  #1 & #2
  \end{tabularx}\par}%
}

\makeatletter
% Make \input{} resilient to different build working directories.
% This is needed because some editors/CI compile from repo root (or a wrapper
% main.tex), while this manuscript lives in new/src/paper/.
\def\input@path{{./}{tables/}{figures/}{src/paper/}{src/paper/tables/}{src/paper/figures/}{new/src/paper/}{new/src/paper/tables/}{new/src/paper/figures/}}
\makeatother

\makeatletter
\IfFileExists{tables/macros_v12.tex}{% macros_v12.tex — auto-generated by scripts/generate_figures_v12.py
\makeatletter
\providecommand{\VNepisodes}{80}
\providecommand{\VQoNoShield}{3270}
\providecommand{\VQoShield}{3218}
\providecommand{\VQoRiskGate}{3225}
\providecommand{\VQoPen}{3269}
\providecommand{\VQoMPC}{3275}
\providecommand{\VQoGenie}{3365}
\providecommand{\VRbNoShield}{14.73}
\providecommand{\VRbShieldLo}{1.20}
\providecommand{\VRbShieldHi}{1.45}
\providecommand{\VRbPen}{14.99}
\providecommand{\VRbMPC}{30.70}
\makeatother
}{\typeout{^^J[WARN] Missing tables/macros_v12.tex --- run generate_figures_v12.py^^J}}
\IfFileExists{tables/macros_equivalence.tex}{% macros_equivalence.tex --- auto-generated by scripts/compute_vmaf_equivalence.py
% Post-hoc TOST equivalence test for paired VMAF differences (VMAF-aware vs. legacy shield).
\makeatletter
\providecommand{\VEqMargin}{2.0}
\providecommand{\VEqDEight}{0.318}
\providecommand{\VEqCILoEight}{0.096}
\providecommand{\VEqCIHiEight}{0.539}
\providecommand{\VEqPEight}{6.0\!\times\!10^{-21}}
\providecommand{\VEqDOne}{0.224}
\providecommand{\VEqCILoOne}{0.027}
\providecommand{\VEqCIHiOne}{0.422}
\providecommand{\VEqPOne}{4.3\!\times\!10^{-25}}
\providecommand{\VEqDFiveGEight}{0.014}
\providecommand{\VEqCILoFiveGEight}{-0.073}
\providecommand{\VEqCIHiFiveGEight}{0.101}
\providecommand{\VEqPFiveGEight}{7.7\!\times\!10^{-53}}
\providecommand{\VEqDFiveGOne}{0.063}
\providecommand{\VEqCILoFiveGOne}{-0.004}
\providecommand{\VEqCIHiFiveGOne}{0.131}
\providecommand{\VEqPFiveGOne}{1.9\!\times\!10^{-60}}
\makeatother
}{\typeout{^^J[WARN] Missing tables/macros_equivalence.tex --- run compute_vmaf_equivalence.py^^J}}
\IfFileExists{tables/macros_latency.tex}{% macros_latency.tex --- auto-generated by scripts/benchmark_shield_latency.py
% Wall-clock latency of safe_adjust_action() (VMAF-aware shield, sigma_soft=0.8, tau=8),
% measured over 2880 calls across the 4-video/20-seed evaluation suite
% on the machine used to run this script (single CPU core, Python, no GPU).
\makeatletter
\providecommand{\VLatMeanUs}{8.8}
\providecommand{\VLatMedianUs}{5.3}
\providecommand{\VLatPNinetyNineUs}{28.6}
\providecommand{\VLatMaxUs}{48.3}
\providecommand{\VLatNcalls}{2880}
\makeatother
}{\typeout{^^J[WARN] Missing tables/macros_latency.tex --- run benchmark_shield_latency.py^^J}}
\IfFileExists{tables/macros_ladder.tex}{% Generated by analyze_vmaf_ladder_perchunk.py --- do not edit by hand.
\providecommand{\LadNChunks}{173}
\providecommand{\LadNPairs}{865}
\providecommand{\LadPctChunksInverted}{25.4}
\providecommand{\LadPctChunksFlat}{74.0}
\providecommand{\LadPctPairsInverted}{5.1}
\providecommand{\LadPctPairsFlat}{20.6}
\providecommand{\LadGapErrMean}{1.91}
\providecommand{\LadGapErrPNinety}{4.48}
\providecommand{\LadGapErrMax}{19.72}
\providecommand{\LadFlatThresh}{1.0}
\providecommand{\LadWidestVideo}{BigBuckBunny}
\providecommand{\LadWidestPair}{750\to1200}
\providecommand{\LadWidestMin}{3.73}
\providecommand{\LadWidestMax}{28.47}
\providecommand{\LadCheapVideo}{BigBuckBunny}
\providecommand{\LadCheapPair}{2850\to6000}
\providecommand{\LadCheapMean}{-0.31}
}{\typeout{^^J[WARN] Missing tables/macros_ladder.tex --- run analyze_vmaf_ladder_perchunk.py^^J}}
\IfFileExists{tables/macros_ladder_v19.tex}{% Session-mean ladder macros (v19) — make_cps_paper_assets.py
\providecommand{\LadNVideos}{12}
\providecommand{\LadPctNonMonotone}{67}
\providecommand{\LadTopGainMean}{5.61}
}{\typeout{^^J[WARN] Missing tables/macros_ladder_v19.tex --- run make_cps_paper_assets.py^^J}}
\IfFileExists{tables/macros_cps.tex}{% Auto-generated by make_cps_paper_assets.py -- DO NOT EDIT BY HAND.
% Certified Perceptual Shield (V18) numbers.
\newcommand{\CPSepisodes}{204}
\newcommand{\CPSepsilon}{1.0}
\newcommand{\CPSalpha}{0.10}
\newcommand{\CPScovTarget}{0.90}
\newcommand{\CPSGreedyRawReb}{85.5}
\newcommand{\CPSGreedyRawVM}{89.1}
\newcommand{\CPSGreedyRawBr}{6000}
\newcommand{\CPSGreedyRawInt}{0.0}
\newcommand{\CPSGreedySafReb}{41.9}
\newcommand{\CPSGreedySafVM}{79.4}
\newcommand{\CPSGreedySafBr}{3131}
\newcommand{\CPSGreedySafCov}{0.919}
\newcommand{\CPSGreedySafInt}{65.8}
\newcommand{\CPSGreedyCertReb}{38.0}
\newcommand{\CPSGreedyCertVM}{79.4}
\newcommand{\CPSGreedyCertBr}{2990}
\newcommand{\CPSGreedyCertCov}{0.919}
\newcommand{\CPSGreedyCertInt}{70.4}
\newcommand{\CPSBbaRawReb}{78.2}
\newcommand{\CPSBbaRawVM}{85.4}
\newcommand{\CPSBbaRawBr}{4662}
\newcommand{\CPSBbaRawInt}{0.0}
\newcommand{\CPSBbaSafReb}{42.1}
\newcommand{\CPSBbaSafVM}{78.4}
\newcommand{\CPSBbaSafBr}{2809}
\newcommand{\CPSBbaSafCov}{0.919}
\newcommand{\CPSBbaSafInt}{58.1}
\newcommand{\CPSBbaCertReb}{38.3}
\newcommand{\CPSBbaCertVM}{78.5}
\newcommand{\CPSBbaCertBr}{2692}
\newcommand{\CPSBbaCertCov}{0.919}
\newcommand{\CPSBbaCertInt}{66.1}
\newcommand{\CPSPenRawReb}{0.1}
\newcommand{\CPSPenRawVM}{64.2}
\newcommand{\CPSPenRawBr}{723}
\newcommand{\CPSPenRawInt}{0.0}
\newcommand{\CPSPenSafReb}{0.1}
\newcommand{\CPSPenSafVM}{63.2}
\newcommand{\CPSPenSafBr}{693}
\newcommand{\CPSPenSafCov}{0.919}
\newcommand{\CPSPenSafInt}{6.7}
\newcommand{\CPSPenCertReb}{0.1}
\newcommand{\CPSPenCertVM}{63.2}
\newcommand{\CPSPenCertBr}{693}
\newcommand{\CPSPenCertCov}{0.919}
\newcommand{\CPSPenCertInt}{6.7}
\newcommand{\CPSBolaRawReb}{73.7}
\newcommand{\CPSBolaRawVM}{80.8}
\newcommand{\CPSBolaRawBr}{4609}
\newcommand{\CPSBolaRawInt}{0.0}
\newcommand{\CPSBolaSafReb}{38.5}
\newcommand{\CPSBolaSafVM}{74.0}
\newcommand{\CPSBolaSafBr}{2753}
\newcommand{\CPSBolaSafCov}{0.919}
\newcommand{\CPSBolaSafInt}{53.6}
\newcommand{\CPSBolaCertReb}{25.0}
\newcommand{\CPSBolaCertVM}{74.2}
\newcommand{\CPSBolaCertBr}{2306}
\newcommand{\CPSBolaCertCov}{0.919}
\newcommand{\CPSBolaCertInt}{69.8}
\newcommand{\CPSMpcRawReb}{62.1}
\newcommand{\CPSMpcRawVM}{73.8}
\newcommand{\CPSMpcRawBr}{3256}
\newcommand{\CPSMpcRawInt}{0.0}
\newcommand{\CPSMpcSafReb}{33.7}
\newcommand{\CPSMpcSafVM}{71.4}
\newcommand{\CPSMpcSafBr}{2371}
\newcommand{\CPSMpcSafCov}{0.919}
\newcommand{\CPSMpcSafInt}{30.4}
\newcommand{\CPSMpcCertReb}{20.4}
\newcommand{\CPSMpcCertVM}{71.5}
\newcommand{\CPSMpcCertBr}{2009}
\newcommand{\CPSMpcCertCov}{0.919}
\newcommand{\CPSMpcCertInt}{41.0}
\newcommand{\CPSGreedyBWcs}{-4.5}
\newcommand{\CPSGreedyRebCs}{-9.3}
\newcommand{\CPSGreedyVMcs}{-0.01}
\newcommand{\CPSGreedyRebCr}{-55.5}
\newcommand{\CPSGreedyCov}{0.919}
\newcommand{\CPSGreedyPBWcs}{8.5\!\times\!10^{-4}}
\newcommand{\CPSGreedyPRebCs}{8.5\!\times\!10^{-4}}
\newcommand{\CPSBbaBWcs}{-4.1}
\newcommand{\CPSBbaRebCs}{-8.9}
\newcommand{\CPSBbaVMcs}{0.17}
\newcommand{\CPSBbaRebCr}{-51.0}
\newcommand{\CPSBbaCov}{0.919}
\newcommand{\CPSBbaPBWcs}{4.2\!\times\!10^{-16}}
\newcommand{\CPSBbaPRebCs}{2.5\!\times\!10^{-4}}
\newcommand{\CPSBolaBWcs}{-16.3}
\newcommand{\CPSBolaRebCs}{-35.0}
\newcommand{\CPSBolaVMcs}{0.14}
\newcommand{\CPSBolaRebCr}{-66.0}
\newcommand{\CPSBolaCov}{0.919}
\newcommand{\CPSBolaPBWcs}{5.0\!\times\!10^{-5}}
\newcommand{\CPSBolaPRebCs}{5.0\!\times\!10^{-5}}
\newcommand{\CPSMpcBWcs}{-15.2}
\newcommand{\CPSMpcRebCs}{-39.2}
\newcommand{\CPSMpcVMcs}{0.13}
\newcommand{\CPSMpcRebCr}{-67.1}
\newcommand{\CPSMpcCov}{0.919}
\newcommand{\CPSMpcPBWcs}{5.0\!\times\!10^{-5}}
\newcommand{\CPSMpcPRebCs}{5.0\!\times\!10^{-5}}
\newcommand{\CPSGreedyImpBWcs}{-14.6}
\newcommand{\CPSGreedyImpRebCs}{-12.6}
\newcommand{\CPSGreedyImpVMcs}{-0.23}
\newcommand{\CPSGreedyImpCov}{0.919}
\newcommand{\CPSBbaImpBWcs}{-12.5}
\newcommand{\CPSBbaImpRebCs}{-13.0}
\newcommand{\CPSBbaImpVMcs}{-0.02}
\newcommand{\CPSBbaImpCov}{0.919}
\newcommand{\CPSbbBWcs}{-0.9}
\newcommand{\CPSbbRebSr}{-98.7}
\newcommand{\CPSPropBWcs}{-4.7}
\newcommand{\CPSPropRebCs}{-20.5}
\newcommand{\CPSPropVMcs}{0.02}
\newcommand{\CPSPropPRebCs}{2.5\!\times\!10^{-3}}
\newcommand{\CPSPropImpBWcs}{-5.7}
\newcommand{\CPSPropImpRebCs}{-27.2}
\newcommand{\CPSPropImpVMcs}{-0.03}
\newcommand{\CPSPropImpPRebCs}{1.1\!\times\!10^{-3}}
\newcommand{\CPScoRegVM}{69.72}
\newcommand{\CPScoRegReb}{0.71}
\newcommand{\CPScoRegBr}{1071}
\newcommand{\CPScoRegCov}{0.919}
\newcommand{\CPScoCpsVM}{72.31}
\newcommand{\CPScoCpsReb}{1.28}
\newcommand{\CPScoCpsBr}{1489}
\newcommand{\CPScoCpsCov}{0.919}
\newcommand{\CPScoDVM}{2.60}
\newcommand{\CPScoDReb}{0.57}
\newcommand{\CPScoWstar}{4.58}
\newcommand{\CPScoPVM}{5.0\!\times\!10^{-5}}
\newcommand{\CPScoPReb}{4.5\!\times\!10^{-3}}
\newcommand{\CPScoDVMlo}{2.30}
\newcommand{\CPScoDVMhi}{2.90}
\newcommand{\CPScoDRebLo}{0.18}
\newcommand{\CPScoDRebHi}{0.97}
\newcommand{\CPScoWstarLo}{2.65}
\newcommand{\CPScoWstarHi}{14.10}
\newcommand{\CPScoWstarMed}{4.57}
\newcommand{\CPScoDVMposPct}{100.0}
\newcommand{\CPScoBootN}{200}
}{\typeout{^^J[WARN] Missing tables/macros_cps.tex --- run make_cps_macros.py^^J}}
\IfFileExists{tables/macros_perchunk.tex}{% Auto-generated per-chunk ablation macros
\newcommand{\PCnChunks}{113}
\newcommand{\PCinvPct}{26.5}
\newcommand{\PCeps}{1.0}
\newcommand{\PCkneePooled}{4.15}
\newcommand{\PCkneePerchunk}{3.94}
\newcommand{\PCbwPooled}{44.6}
\newcommand{\PCbwPerchunk}{45.8}
\newcommand{\PCcostPerchunk}{0.32}
\newcommand{\PCcostPooled}{0.33}
\newcommand{\PCviolPooled}{11.5}
\newcommand{\PCdisagree}{33.6}
}{\typeout{^^J[WARN] Missing tables/macros_perchunk.tex --- run ablation_perchunk_ladder.py^^J}}
\IfFileExists{tables/macros_ablation.tex}{% Auto-generated eps/alpha ablation macros
\newcommand{\AblEpsLo}{0.5}
\newcommand{\AblEpsHi}{4.0}
\newcommand{\AblBwLo}{8.5}
\newcommand{\AblBwHi}{38.4}
\newcommand{\AblVmMaxAbs}{0.42}
\newcommand{\AblCovMin}{0.854}
\newcommand{\AblCovMax}{0.945}
}{\typeout{^^J[WARN] Missing tables/macros_ablation.tex --- run ablation_eps_alpha.py^^J}}
\makeatother

% CPS fallbacks (overridden by macros_cps.tex when present)
\providecommand{\CPSepisodes}{204}
\providecommand{\CPSepsilon}{1.0}
\providecommand{\CPSalpha}{0.10}
\providecommand{\CPScovTarget}{0.90}
\providecommand{\CPSGreedyBWcs}{-4.5}
\providecommand{\CPSGreedyRebCs}{-9.3}
\providecommand{\CPSGreedyVMcs}{-0.01}
\providecommand{\CPSGreedyRebCr}{-55.5}
\providecommand{\CPSGreedyCov}{0.919}
\providecommand{\CPSGreedyPBWcs}{8.5\!\times\!10^{-4}}
\providecommand{\CPSGreedyPRebCs}{8.5\!\times\!10^{-4}}
\providecommand{\CPSBbaBWcs}{-4.1}
\providecommand{\CPSBbaRebCs}{-8.9}
\providecommand{\CPSBbaVMcs}{0.17}
\providecommand{\CPSBbaRebCr}{-51.0}
\providecommand{\CPSBbaCov}{0.919}
\providecommand{\CPSBbaPBWcs}{4.2\!\times\!10^{-16}}
\providecommand{\CPSBbaPRebCs}{2.5\!\times\!10^{-4}}
\providecommand{\CPSGreedyImpBWcs}{-14.6}
\providecommand{\CPSGreedyImpRebCs}{-12.6}
\providecommand{\CPSGreedyImpVMcs}{-0.23}
\providecommand{\CPSBbaImpBWcs}{-12.5}
\providecommand{\CPSBbaImpRebCs}{-13.0}
\providecommand{\CPSBbaImpVMcs}{-0.02}
\providecommand{\CPSbbBWcs}{-0.9}
\providecommand{\CPSbbRebSr}{-98.7}
\providecommand{\CPSbbCov}{0.890}
\providecommand{\CPSBolaBWcs}{-4.1}
\providecommand{\CPSBolaRebCs}{-8.9}
\providecommand{\CPSBolaVMcs}{0.18}
\providecommand{\CPSMpcBWcs}{-3.8}
\providecommand{\CPSMpcRebCs}{-9.2}
\providecommand{\CPSMpcVMcs}{0.02}
\providecommand{\CPSPropBWcs}{-35.3}
\providecommand{\CPSPropRebCs}{-79.3}
\providecommand{\CPSPropVMcs}{-0.33}
\providecommand{\CPSPropPRebCs}{6.2\!\times\!10^{-18}}
\providecommand{\CPSPropImpBWcs}{-21.9}
\providecommand{\CPSPropImpRebCs}{-29.6}
\providecommand{\CPSPropImpVMcs}{-0.27}
\providecommand{\CPSPropImpPRebCs}{3.2\!\times\!10^{-14}}
\providecommand{\CPScoRegVM}{72.80}
\providecommand{\CPScoRegReb}{0.31}
\providecommand{\CPScoRegBr}{801}
\providecommand{\CPScoRegCov}{0.919}
\providecommand{\CPScoCpsVM}{75.73}
\providecommand{\CPScoCpsReb}{0.86}
\providecommand{\CPScoCpsBr}{1463}
\providecommand{\CPScoCpsCov}{0.919}
\providecommand{\CPScoDVM}{2.93}
\providecommand{\CPScoDReb}{0.55}
\providecommand{\CPScoWstar}{5.31}
\providecommand{\CPScoPVM}{5.0\!\times\!10^{-5}}
\providecommand{\CPScoPReb}{5.0\!\times\!10^{-5}}
\providecommand{\CPScoDVMlo}{2.61}
\providecommand{\CPScoDVMhi}{3.25}
\providecommand{\CPScoDRebLo}{0.29}
\providecommand{\CPScoDRebHi}{0.85}
\providecommand{\CPScoWstarLo}{3.40}
\providecommand{\CPScoWstarHi}{10.27}
\providecommand{\CPScoWstarMed}{5.37}
\providecommand{\CPScoDVMposPct}{100.0}
\providecommand{\CPScoBootN}{204}

% -------------------------------------------------------------------------
% Fallback macros (compile-safety)
% If this file is compiled from a different working directory, the relative
% path above may fail and macros will be undefined. These placeholders ensure
% the paper still compiles; they will be overridden when macros_v12.tex loads.
% -------------------------------------------------------------------------
\providecommand{\VNepisodes}{80}
\providecommand{\VShieldVariants}{47}
\providecommand{\VShieldBehaviors}{11}
\providecommand{\VPairedDQoEvsLeg}{+42.13}
\providecommand{\VPairedPQoEvsLeg}{5.0\!\times\!10^{-3}}
\providecommand{\VPairedDRbvsLeg}{-0.55}
\providecommand{\VPairedPRbvsLeg}{2.9\!\times\!10^{-3}}
\providecommand{\VPairedDVMvsLeg}{+0.22}
\providecommand{\VPairedDQoEvsLegStr}{+64.82}
\providecommand{\VPairedDRbvsLegStr}{-0.69}
\providecommand{\VPairedDVMvsLegStr}{+0.32}
\providecommand{\VPairedPVMvsLeg}{0.056}
\providecommand{\VPairedPQoEvsLegStr}{1.4\!\times\!10^{-3}}
\providecommand{\VPairedPRbvsLegStr}{2.7\!\times\!10^{-3}}
\providecommand{\VPairedPVMvsLegStr}{4.1\!\times\!10^{-2}}
\providecommand{\VPairedDRbvsOff}{-5.92}
\providecommand{\VPairedPRbvsOff}{1.4\!\times\!10^{-10}}
\providecommand{\VRbReductionVsPen}{9.87}
\providecommand{\VRbReductionVsMPC}{25.58}
\providecommand{\VRbPen}{14.99}
\providecommand{\VRbMPC}{30.70}
\providecommand{\VRbVmAwareOne}{5.12}
\providecommand{\VRbOff}{11.04}
\providecommand{\VFiveGRbShieldOff}{70.26}
\providecommand{\VFiveGRbShieldLegacy}{21.91}
\providecommand{\VFiveGPairedDQoE}{+39.85}
\providecommand{\VFiveGPairedPQoE}{3.9\!\times\!10^{-5}}
\providecommand{\VFiveGPairedDRb}{-2.38}
\providecommand{\VFiveGPairedPRb}{4.4\!\times\!10^{-6}}
\providecommand{\VFiveGPairedDVm}{+0.01}
\providecommand{\VFiveGPairedPVm}{0.715}
\providecommand{\VFiveGPairedDQoEten}{+34.60}
\providecommand{\VFiveGPairedPQoEten}{8.1\!\times\!10^{-7}}
\providecommand{\VFiveGPairedDRbten}{-2.27}
\providecommand{\VFiveGPairedPRbten}{2.8\!\times\!10^{-6}}
\providecommand{\VFiveGPairedPVmten}{0.422}

% --- Post-hoc TOST equivalence test for paired VMAF differences (see
% scripts/compute_vmaf_equivalence.py; macros_equivalence.tex overrides these). ---
\providecommand{\VEqMargin}{2.0}
\providecommand{\VEqDEight}{0.318}
\providecommand{\VEqCILoEight}{0.096}
\providecommand{\VEqCIHiEight}{0.539}
\providecommand{\VEqPEight}{6.0\!\times\!10^{-21}}
\providecommand{\VEqDOne}{0.224}
\providecommand{\VEqCILoOne}{0.027}
\providecommand{\VEqCIHiOne}{0.422}
\providecommand{\VEqPOne}{4.3\!\times\!10^{-25}}
\providecommand{\VEqDFiveGEight}{0.014}
\providecommand{\VEqCILoFiveGEight}{-0.073}
\providecommand{\VEqCIHiFiveGEight}{0.101}
\providecommand{\VEqPFiveGEight}{7.7\!\times\!10^{-53}}
\providecommand{\VEqDFiveGOne}{0.063}
\providecommand{\VEqCILoFiveGOne}{-0.004}
\providecommand{\VEqCIHiFiveGOne}{0.131}
\providecommand{\VEqPFiveGOne}{1.9\!\times\!10^{-60}}

% --- Shield latency micro-benchmark (see scripts/benchmark_shield_latency.py;
% macros_latency.tex overrides these). ---
\providecommand{\VLatMeanUs}{8.8}
\providecommand{\VLatMedianUs}{5.3}
\providecommand{\VLatPNinetyNineUs}{28.6}
\providecommand{\VLatMaxUs}{48.3}
\providecommand{\VLatNcalls}{2880}

% --- Pareto family table macros (from v123 shield sweep) ---
% Legacy (shield_legacy, ρ_cat=2)
\providecommand{\VQoLegacy}{2518.87}
\providecommand{\VRbLegacy}{5.67}
\providecommand{\VVmLegacy}{65.82}
\providecommand{\VIntervLegacy}{3.23}
% VMAF-aware τ=0.8
\providecommand{\VQoVmAwareEpt}{2583.70}
\providecommand{\VRbVmAwareEpt}{4.98}
\providecommand{\VVmVmAwareEpt}{66.13}
\providecommand{\VIntervVmAwareEpt}{2.79}
% VMAF-aware τ=1.0
\providecommand{\VQoVmAwareOne}{2561.00}
\providecommand{\VRbVmAwareOne}{5.12}
\providecommand{\VVmVmAwareOne}{66.04}
\providecommand{\VIntervVmAwareOne}{2.71}
% VMAF-aware τ=1.2
\providecommand{\VQoVmAwareTwl}{2546.86}
\providecommand{\VRbVmAwareTwl}{5.03}
\providecommand{\VVmVmAwareTwl}{65.89}
% VMAF-aware τ=1.5
\providecommand{\VQoVmAwareFvt}{2532.23}
\providecommand{\VRbVmAwareFvt}{5.21}
\providecommand{\VVmVmAwareFvt}{65.86}
% Threshold only (no VMAF FB)
\providecommand{\VQoThreshThree}{2519.40}
\providecommand{\VRbThreshThree}{8.94}
\providecommand{\VQoThreshFour}{2530.78}
\providecommand{\VRbThreshFour}{9.96}
\providecommand{\VQoThreshFive}{2557.17}
\providecommand{\VRbThreshFive}{10.49}
% Threshold + VMAF FB
\providecommand{\VQoThreshThreeVmafFB}{2539.11}
\providecommand{\VRbThreshThreeVmafFB}{8.45}
\providecommand{\VVmThreshThreeVmafFB}{66.30}
\providecommand{\VQoThreshFourVmafFB}{2526.21}
\providecommand{\VRbThreshFourVmafFB}{9.95}
\providecommand{\VVmThreshFourVmafFB}{66.36}
% Shield off
\providecommand{\VQoOff}{2563.72}
\providecommand{\VVmOff}{66.55}

% Target journal (Elsevier metadata; Computer Networks guide for authors).
\journal{Computer Networks}

\begin{document}

\begin{frontmatter}

\title{Certified Perceptual Shielding for Adaptive Bitrate Streaming}

\author{\mbox{Saeed Zarbi}$^{a}$}
\ead{szarbi@tabrizu.ac.ir}
\author{\mbox{Leili Farzinvash}$^{a}$}
\ead{l.farzinvash@tabrizu.ac.ir}
\author{\mbox{Pedram Salehpour}$^{a}$\corref{c1}}
\ead{psalehpour@tabrizu.ac.ir}
\cortext[c1]{Corresponding author: psalehpour@tabrizu.ac.ir}

\address{$^{a}$Department of Computer Engineering, Faculty of Electrical and Computer Engineering, University of Tabriz, Tabriz, Iran}

\begin{abstract}
Client-side adaptive bitrate (ABR) controllers---whether heuristic, model-based, or learned---share a per-chunk limitation.
A request can become unsafe when throughput fluctuates, or wasteful once the rate--quality ladder has saturated in Video Multimethod Assessment Fusion (VMAF).
We addressed both failure modes with a \emph{Certified Perceptual Shield} (CPS).
CPS is a model-agnostic runtime projection built from three components.
\emph{VMAF-knee bandwidth banking} lowers even feasible proposals to the smallest rung whose VMAF lies within a perceptual budget \(\varepsilon\) of the proposed rung.
The mechanism exploits saturation to bank buffer headroom with negligible quality loss.
An online \emph{conformal} lower bound on throughput is calibrated with a finite-sample \((n{+}1)\) rank correction.
Under exchangeability, this bound yields distribution-free coverage \(1{-}\alpha\).
\emph{Certified feasibility} then rejects any download whose time under that bound would exceed the buffer margin.
The result is a per-chunk \emph{probabilistic} stall certificate under the conformal assumption.
The certificate is a coverage guarantee, not a promise of zero stalls.
We evaluated CPS on \(\CPSepisodes\) paired synthetic \mbox{5G} episodes over twelve reference titles.
The titles were chosen for \emph{ladder diversity}, from saturation-rich animation to steep natural motion.
Against a conformal safety-only baseline, banking significantly reduced delivered bitrate by about 4.5--3.8\% and rebuffering by about 9.3--9.2\% on every deterministic host we tested (greedy through RobustMPC).
Mean VMAF differences lie within \(\pm\CPSepsilon\) (TOST).
Empirical coverage was \(\CPSGreedyCov\) against a \(\CPScovTarget\) target, at or above the nominal level.
Because effect sizes scale with top-rung saturation, the diversified roster yielded conservative but credible savings.
A risk-aware, predictive variant increased banking ahead of forecast throughput dips; on greedy it reduced bitrate by 14.6\% and rebuffering by 12.6\%.
On source-disjoint real-measured broadband traces (FCC and Norway/Nornet), banking gains narrowed to about 0.9\% bitrate reduction (certified vs.\ safety).
Conformal safety still reduced stalls by about 98.7\% vs.\ raw.
The payoff is therefore regime-dependent.
Optional PPO co-training with CPS was reported only as a supporting, single-seed study in the evaluation---not as a headline claim.
\end{abstract}

\begin{keyword}
Adaptive Bitrate Streaming\sep Runtime Shielding\sep VMAF\sep Conformal Prediction\sep Certified Safety\sep Video QoE
\end{keyword}

\end{frontmatter}

% =========================================
\section{Introduction}
\label{sec:intro}

Client-side adaptive bitrate (ABR) policies optimize average quality of experience (QoE) through buffer heuristics~\cite{huang2014buffer}, model-predictive control~\cite{yin2015control}, and deep reinforcement learning (DRL)~\cite{mao2017neural,schulman2017proximal}.
Session averages, however, can obscure transient risks associated with individual chunks.
A single segment can still stall under volatile sub-6\,GHz and \mbox{5G/mmWave} throughput~\cite{narayanan2021variegated,cisco2023,seufert2015survey}, whereas a comfortable buffer invites the opposite inefficiency.
Commercial rate ladders often saturate in Video Multimethod Assessment Fusion (VMAF)~\cite{li2016vmaf}: further bitrate increases at the top often provide only marginal perceptual gains.
Requesting those rungs when the buffer is healthy consumes bandwidth that could cushion the next dip.

A runtime shield sits between the base policy and the player, projecting each proposal before download~\cite{alshiekh2018shielding}.
The standard repair decrements the representation index until a buffer-feasibility test passes.
That repair is content-blind: it ignores VMAF saturation, never banks bandwidth proactively, and reduces throughput pessimism to a hand-tuned factor with no coverage guarantee.
Deployed ladders carry exploitable structure---rate--quality saturation near the top.
That structure supports perceptually negligible bitrate reductions while freeing additional buffer capacity.

We introduced a \emph{Certified Perceptual Shield} (CPS) that targets this structure while stating a formal coverage guarantee:
\begin{enumerate}
  \item \textbf{VMAF-knee bandwidth banking.} Even a feasible action is lowered to the smallest rung \(j\) whose VMAF lies within a budget \(\varepsilon\) of the proposal. The downshift is perceptually bounded, and reclaimed bytes bank as buffer for later chunks.
  \item \textbf{Conformal throughput lower bound.} Download-time tests use an online split-conformal lower bound on throughput with target coverage \(1{-}\alpha\). The bound is calibrated from residual ratios with a finite-sample \((n{+}1)\) rank correction~\cite{vovk2005algorithmic,angelopoulos2023gentle}.
  \item \textbf{Certified feasibility.} After banking, CPS admits no rung whose download time under the lower bound exceeds the buffer minus a safety margin. This rule yields a per-chunk probabilistic stall certificate with coverage \(1{-}\alpha\) under residual exchangeability, rather than a deterministic no-stall guarantee (Proposition~\ref{prop:feas}).
\end{enumerate}
Prior VMAF-aware ranking shields provably collapse to highest-feasible-index projection on monotone ladders (Proposition~\ref{prop:collapse}).
A perceptual-loss budget therefore cannot constrain choices when VMAF is nondecreasing in the rung index.
The wrapper is host-agnostic: the same projection attaches to bitrate-greedy, buffer-based (BBA), and learned PPO controllers alike (optional co-design in Section~\ref{sec:eval:codesign}).
To separate the two mechanisms, we ran three paired arms on identical episode seeds: \textsc{Raw} (unshielded), \textsc{Safety} (conformal feasibility only), and \textsc{Certified} (feasibility plus banking).
Banking was then isolated to the \textsc{Certified}--\textsc{Safety} contrast, and safety to \textsc{Safety}/\textsc{Certified} vs.\ \textsc{Raw}.

We evaluated these claims on \(\CPSepisodes\) paired synthetic \mbox{5G} episodes over twelve titles selected for ladder diversity, with \(\varepsilon{=}\CPSepsilon\) and \(\alpha{=}\CPSalpha\).
Certified banking against safety-only significantly reduced mean bitrate and rebuffering on greedy, BBA, BOLA, and RobustMPC (\tabref{tab:cps:full}).
For greedy, banking reduced bitrate by 4.5\% and rebuffering by 9.3\% (paired Wilcoxon \(p{=}\CPSGreedyPRebCs\)).
Magnitudes clustered near 4--5\% bitrate and about 9\% rebuffering because steep, non-saturated ladders left fewer near-lossless downshifts, yet the direction held everywhere.
Mean VMAF differences remained within \(\pm\varepsilon\) by TOST equivalence.
Empirical conformal coverage was \(\CPSGreedyCov\) against a \(\CPScovTarget\) target.
A risk-aware, predictive extension widened the perceptual budget when the buffer---or a short look-ahead under a throughput floor---entered a danger zone.
On greedy, this variant reduced bitrate by 14.6\% and rebuffering by 12.6\%.
On source-disjoint real-measured broadband traces (FCC and Norway/Nornet), banking gains narrowed to about 0.9\% bitrate reduction (certified vs.\ safety).
Conformal safety still reduced stalls by about 98.7\% vs.\ raw greedy.
These traces therefore indicated when perceptual banking paid off.
Optional PPO co-training with CPS was deferred to a supporting, single-seed study (Section~\ref{sec:eval:codesign}).

The main contributions are:
\begin{itemize}
  \item \textbf{Certified Perceptual Shield.} A model-agnostic runtime projection that combines VMAF-knee bandwidth banking, a conformal throughput lower bound with finite-sample coverage, and certified buffer feasibility (Section~\ref{sec:method}).

  \item \textbf{Isolated banking evaluation.} Paired \textsc{Raw}/\textsc{Safety}/\textsc{Certified} rollouts on greedy, BBA, BOLA, RobustMPC, and PPO hosts. We reported Wilcoxon tests on bitrate and rebuffering, and TOST non-inferiority on VMAF within \(\pm\varepsilon\) (Section~\ref{sec:eval}).

  \item \textbf{Risk-aware, predictive banking.} A perceptual budget conditioned on the buffer and a short look-ahead that pre-banks before forecast dips. It improves the trade-off between bandwidth and stalls under the same \(\varepsilon\) ceiling at comfortable buffers (Section~\ref{sec:method:predictive}).

  \item \textbf{Regime characterization.} Banking gains concentrated where top rungs were often feasible (synthetic \mbox{5G} with a moderate buffer). Real-trace broadband stress was dominated by safety projection (Section~\ref{sec:eval}).

  \item \textbf{Policy--shield co-design (supporting).} An optional single-seed PPO study with CPS in the training loop; we reported the QoE crossover weight \(w^{\star}\) but did not treat this arm as a headline claim (Section~\ref{sec:eval:codesign}).
\end{itemize}

The remainder of this paper is organized as follows.
Section~\ref{sec:related} reviews related work, and Section~\ref{sec:system} formalizes the streaming model.
Section~\ref{sec:method} presents CPS and optional co-training.
Section~\ref{sec:eval} reports the evaluation, Section~\ref{sec:limits} discusses limitations, and Section~\ref{sec:conclusion} concludes.

% =========================================
\section{Related work}
\label{sec:related}

Three areas bear on client-side representation selection: model-based ABR, learned ABR, and constrained control with runtime shielding.
Our focus is the control policy and its runtime action projection; throughput prediction lies outside this scope.
Predictors such as CS2P~\cite{sun2016cs2p} play a complementary role: they feed inputs to model-based or learned controllers~\cite{yin2015control,mao2017neural}.
ABR surveys~\cite{bentaleb2018survey,seufert2015survey,peroni2025survey} and \mbox{5G/mmWave} measurements~\cite{narayanan2021variegated,fezeu2024midband} converge on one empirical lesson.
Control must remain robust when throughput fluctuates sharply.

\subsection{Heuristic and model-based ABR}

MPEG Dynamic Adaptive Streaming over HTTP (DASH) separates packaging from transport~\cite{sodagar2011mpeg}.
Clients execute local adaptation policies based on ladder metadata and observed download conditions.
Buffer-threshold heuristics such as BBA~\cite{huang2014buffer} remain simple and interpretable.
BOLA~\cite{spiteri2018bola}, by contrast, derived its decisions from a Lyapunov-based utility formulation.
RobustMPC~\cite{yin2015control} planned ahead.
It selects representations by solving a short-horizon model predictive control (MPC) problem under estimated future throughput.
The payoff is interpretable thresholds and horizons.
The cost is fragility: performance depends on manual calibration and on prediction accuracy that varies across networks and devices~\cite{narayanan2021variegated,cisco2023}.
Moreover, conventional feasibility rules rarely exploit VMAF saturation to bank bandwidth in advance~\cite{li2016vmaf}.
Other related efforts target different stages of the delivery pipeline.
SODA~\cite{chen2024soda} pursued smoothness-aware online control, while edge assistance~\cite{mehrabi2019edge} and Oboe-style automatic tuning~\cite{akhtar2018oboe} addressed separate parts of the system.

\subsection{Deep reinforcement learning for ABR}

Pensieve~\cite{mao2017neural} established the template: a trace-driven DRL formulation for ABR.
D-DASH~\cite{gadaleta2017ddash} and more recent PPO-based controllers~\cite{naresh2023ppoabr} followed this line.
Their neural policies combine throughput histories, buffer state, and manifest features~\cite{sutton2018reinforcement}.
Later systems extended the Pensieve trace-driven DRL template.
Comyco~\cite{huang2019comyco} trained by imitation against a VMAF-based QoE objective and showed that perceptual signals helped inside the learned policy.
Loki~\cite{zhang2021loki} combined learned features with a rule-based model and reduced long-tail failures in production; its rule module soft-guards extremes but lacks formal feasibility certification.
Fugu~\cite{yan2019learning} deployed learned control in live networks and exposed the distribution shift between simulation and deployment, which underscored the need for runtime safeguards beyond expected-return training.
Taken together, these systems motivate constrained learning and decision-time enforcement.
An expected-return objective alone encodes no explicit per-episode stall or smoothness budget~\cite{hu2023corel,gu2024safereview}.
Explicit constraints do not fully close this gap.
A controller trained under a CMDP retains guarantees only over the sampled trajectory distribution.
An unseen blockage pattern can still induce an unsafe proposal at a single chunk boundary.
This gap argues for splitting the stack into a learned performance layer and a runtime enforcement layer.
Additional training cannot substitute for runtime validation.

\subsection{Safe/constrained RL, shields, and the content-blindness gap}

CMDPs and Lagrangian primal--dual methods optimize expected return under explicit constraints~\cite{altman1999constrained,paternain2019dual,achiam2017constrained}.
Safe-RL surveys catalogued related families: shielding, Lyapunov, and risk-sensitive schemes~\cite{gu2024safereview,chow2018lyapunov}.
COREL~\cite{hu2023corel} is closest to our training formulation.
It cast ABR over commercial \mbox{5G/mmWave} traces as a CMDP with a stall-time constraint.
The CMDP is solved by a primal--dual actor--critic.
That constraint, however, operated within training; it is not a standalone runtime projection with auditable per-decision intervention logging.
Runtime shielding takes the opposite approach: it repairs a proposed action before execution and can log each intervention~\cite{alshiekh2018shielding}.
Surveys split the family into \emph{pre-shielding}, which prunes the action set, and \emph{post-shielding}, which overwrites an unsafe proposal~\cite{konighofer2025shields}.
By that taxonomy, CPS is a post-shield.
Its safety set is drawn by a conformal feasibility test rather than a temporal-logic model.

This distinction shaped our design: we separated an optional learned performance layer from a deterministic runtime CPS.
Learned policies such as Lagrangian PPO can pursue expected stall and smoothness budgets across training rollouts~\cite{paternain2019dual,hu2023corel}.
They cannot certify per-chunk feasibility under conformal throughput bounds.
CPS targets exactly this gap.
Before every download it (i)~banks bandwidth on VMAF-saturated rungs within a perceptual budget \(\varepsilon\), and (ii)~projects onto actions feasible under a distribution-free throughput lower bound with coverage \(1{-}\alpha\)~\cite{vovk2005algorithmic,angelopoulos2023gentle,alshiekh2018shielding}.
Its use of perceptual data also differs from Comyco, where VMAF enters only the learned objective~\cite{huang2019comyco}.
CPS reads the ladder at decision time for proactive banking, rather than only emergency index repair.
\tabref{tab:rw_safety_vmaf} compares client-side ABR methods by training constraints, runtime guard, conformal bounds, and VMAF-knee banking.

\begin{cps@table}
    \caption{Positioning of client-side ABR methods by training constraints, deterministic runtime projection, conformal throughput bounds, and VMAF-knee banking.}
    \label{tab:rw_safety_vmaf}
    \centering
    \evaltable{%
    \begin{tabularx}{\linewidth}{@{}>{\raggedright\arraybackslash}X l >{\raggedright\arraybackslash}X cccc@{}}
    \toprule
    \textbf{Work} & \textbf{Year} & \textbf{Paradigm} & \textbf{Constr.} & \textbf{Guard} & \textbf{Conf.} & \textbf{Bank} \\
    \midrule
    BBA~\cite{huang2014buffer} & 2014 & Heuristic & N & N & N & N \\
    BOLA~\cite{spiteri2018bola} & 2018 & Lyapunov util. & Imp. & N & N & N \\
    RobustMPC~\cite{yin2015control} & 2015 & MPC/lookahead & N & N & N & N \\
    Pensieve~\cite{mao2017neural} & 2017 & DRL & N & N & N & N \\
    Comyco~\cite{huang2019comyco} & 2019 & Quality-aware IL & N & N & N & N \\
    Loki~\cite{zhang2021loki} & 2021 & Rule--RL fusion & N & \(\approx\) & N & N \\
    COREL~\cite{hu2023corel} & 2023 & Constr.\ DRL & Y & N & N & N \\
    SODA~\cite{chen2024soda} & 2024 & Online control & N & N & N & N \\
    \textbf{Ours (CPS)} & --- & Shield (\(\pm\)RL) & \(\approx\) & \textbf{Y} & \textbf{Y} & \textbf{Y} \\
    \bottomrule
    \end{tabularx}%
    }
    \tablenotecols{%
    \textbf{Constr.:} explicit stall budget in training (Y), implicit (Imp.), optional/partial (\(\approx\)), or none (N).
    \textbf{Guard:} auditable action projection before download~\cite{alshiekh2018shielding}.}{%
    \textbf{Conf.:} conformal throughput lower bound with stated coverage~\cite{vovk2005algorithmic}.
    \textbf{Bank:} VMAF-knee bandwidth banking within a perceptual budget \(\varepsilon\).}%
\end{cps@table}

% =========================================
\section{System model and problem formulation}
\label{sec:system}

This section formalized client-side ABR at the chunk timescale.
A base policy proposes a ladder index; CPS applies a deterministic projection before download; the player then updates the buffer under realized throughput.
The subsections below define the architecture, the state--action--reward tuple, the ladder saturation that banking exploits, and the download and buffer dynamics.
CPS internals---the conformal bound, safety set, and knee rule---are deferred to Section~\ref{sec:method} so \(\underline{C}_k\) is defined before use.
\tabref{tab:notation} collected the symbols used throughout.

\begin{cps@table}
    \caption{Summary of notation.}
    \label{tab:notation}
    \centering
    \evaltable{%
    \begin{tabular}{@{}ll@{}}
    \toprule
    \textbf{Symbol} & \textbf{Meaning} \\
    \midrule
    \(k\), \(T_{\mathrm{c}}\), \(K_e\) & chunk index, chunk duration, chunks per episode \\
    \(L\), \(\mathcal{B}{=}\{0,\dots,L{-}1\}\) & number of ladder rungs, rung index set \\
    \(a^{\mathrm{raw}}_k\), \(a^{\mathrm{exec}}_k\), \(\iota_k\) & proposed rung, executed rung, intervention bit \\
    \(B_k\), \(m\), \(B_{\mathrm{crit}}\) & buffer level, safety margin, critical buffer \\
    \(V_k(j)\), \(\sigma_k(j)\) & VMAF and size (bits) of rung \(j\) at chunk \(k\) \\
    \(\varepsilon\), \(\varepsilon_{\mathrm{eff}}\), \(\varepsilon_{\mathrm{risk}}\) & perceptual budget, effective/risk-mode budget \\
    \(\hat{C}_k\), \(C^{\mathrm{act}}_k\), \(\underline{C}_k\) & throughput predictor, realized, conformal lower bound \\
    \(\alpha\), \(1{-}\alpha\) & conformal miscoverage, target coverage \\
    \(r{=}C^{\mathrm{act}}/\hat{C}\), \(W\), \(\kappa\) & residual ratio, residual window, predictor horizon \\
    \(d_k(j)\), \(H\), \(\varepsilon_0\) & download time of rung \(j\), look-ahead, numerical guard \\
    \(\mu\), \(\beta\) & logged-QoE smoothness and stall weights (\eqref{eq:qoe-rep}) \\
    \(w\), \(w^{\star}\) & QoE stall weight (sweep), mean-QoE crossover weight \\
    \bottomrule
    \end{tabular}%
    }
\end{cps@table}

\subsection{Streaming architecture and control loop}

We worked in the standard Hypertext Transfer Protocol (HTTP) adaptive-streaming abstraction~\cite{yin2015control,mao2017neural,sodagar2011mpeg}.
A video origin or content delivery network (CDN) published \(L\) encoded representations.
We indexed them by \(\mathcal{B}=\{0,\ldots,L{-}1\}\); the DASH/HTTP Live Streaming (HLS) manifest carried the matching per-chunk sizes and perceptual metadata.
At decision step \(k\), the client fetches chunk \(k\) while holding a playout buffer \(B_k\).
The effective trace throughput \(\hat{C}_k\) is exogenous to the controller. This assumption reflects deployment conditions rather than merely simplifying the model. On sub-6\,GHz and \mbox{5G/mmWave} access, throughput can change abruptly because of blockage, line-of-sight (LOS) and non-line-of-sight (NLOS) transitions, scheduling, and mobility~\cite{narayanan2021variegated,cisco2023}.

A base policy \(\pi\) (heuristic or learned) maps an augmented state \(s_k\) to a proposed action \(a^{\mathrm{raw}}_k\in\mathcal{B}\).
Before chunk \(k\) is downloaded, the Certified Perceptual Shield \(\mathcal{S}\) applies the deterministic projection
\begin{equation}
\label{eq:shield-map}
    \big(a^{\mathrm{exec}}_k,\iota_k\big)=\mathcal{S}\!\left(s_k,a^{\mathrm{raw}}_k\right),\qquad a^{\mathrm{exec}}_k\in\mathcal{B},\ \ \iota_k\in\{0,1\},
\end{equation}
The projection returns the executed representation index \(a^{\mathrm{exec}}_k\) and an audit bit \(\iota_k\), the latter indicating whether the projection altered the proposal.
Only \(a^{\mathrm{exec}}_k\) feeds the chunk download and buffer dynamics~\eqref{eq:buffer}; \figref{fig:arch} traces the \(a^{\mathrm{raw}}_k\!\rightarrow\!\mathcal{S}\!\rightarrow\!a^{\mathrm{exec}}_k\) path.

\begin{figure}[!htbp]
    \centering
    \adjustbox{width=\linewidth}{% Figure 1: compact architecture (natural width ~\linewidth, minimal scaling).
\tikzset{
    archbox/.style={
        draw, rounded corners=3pt, align=center,
        inner sep=4pt, font=\small,
        minimum height=0.95cm
    },
    envbox/.style={archbox, fill=gray!10, draw=gray!55, text width=2.0cm},
    statebox/.style={archbox, fill=orange!14, draw=orange!55!black, text width=2.5cm, minimum height=1.25cm},
    actorbox/.style={archbox, fill=blue!10, draw=blue!55!black, text width=2.2cm},
    shieldbox/.style={archbox, fill=green!12, draw=green!45!black, text width=2.4cm, minimum height=1.15cm},
    clientbox/.style={archbox, fill=red!8, draw=red!45!black, text width=1.9cm},
    criticbox/.style={archbox, fill=blue!6, draw=blue!45!black, dashed, text width=2.2cm},
    groupbox/.style={draw=black!30, dashed, rounded corners=4pt, inner xsep=7pt, inner ysep=9pt},
    infer/.style={->, >=Latex, line width=1.05pt, black!80},
    inferbold/.style={->, >=Latex, line width=1.25pt, black},
    train/.style={->, >=Latex, line width=0.8pt, densely dashed, blue!55!black},
    feedback/.style={->, >=Latex, line width=1.05pt, green!45!black},
    edgelabel/.style={font=\footnotesize, fill=white, inner sep=1.5pt},
}

\begin{tikzpicture}[font=\small, >=Latex]

    % --- Environment (left) ---
    \node[envbox] (network) at (0, 1.35) {Network \(C_k\)};
    \node[envbox] (manifest) at (0, 0.45) {Manifest/VMAF};
    \node[envbox] (buffer) at (0, -0.45) {Buffer \(B_k\)};

    % --- Core pipeline ---
    \node[statebox] (state) at (3.3, 0.45)
        {\textbf{State} \(s_k\)\\[-1pt]{\scriptsize throughput, buffer,}\\[-2pt]{\scriptsize prev.\ action, lookahead}};

    \node[actorbox] (actor) at (6.4, 0.45)
        {\textbf{Actor} \(\pi_\theta\)\\[-1pt]{\scriptsize \(a^{\mathrm{raw}}_k\)}};

    \node[shieldbox] (shield) at (9.6, 0.45)
        {\textbf{VMAF-aware Shield}\\[-1pt]{\scriptsize \(\mathcal{S}(s,a^{\mathrm{raw}})\)}};

    \node[clientbox] (client) at (12.4, 0.45) {Client\\Player};

    \node[criticbox] (critic) at (6.4, -1.55)
        {\textbf{Critic} \(V_\phi(s)\)\\[-1pt]{\scriptsize training only}};

    % --- Environment to state ---
    \draw[infer] (network.east) -- ++(0.35,0) |- ([yshift=4pt]state.west);
    \draw[infer] (manifest.east) -- (state.west);
    \draw[infer] (buffer.east) -- ++(0.35,0) |- ([yshift=-4pt]state.west);

    % --- Main inference path (labels above) ---
    \draw[inferbold] (state.east) -- (actor.west);
    \node[edgelabel, above=2pt] at ($(state.east)!0.5!(actor.west)$) {\(s_k\)};

    \draw[inferbold] (actor.east) -- (shield.west);
    \node[edgelabel, above=2pt] at ($(actor.east)!0.5!(shield.west)$) {\(a^{\mathrm{raw}}_k\)};

    \draw[feedback] (shield.east) -- (client.west);
    \node[edgelabel, above=2pt] at ($(shield.east)!0.5!(client.west)$) {\(a^{\mathrm{exec}}_k\)};

    % --- State to shield (lower corridor) ---
    \draw[infer]
        ([yshift=-4pt]state.east) -- ++(0.25,0)
        |- ($(state.east)!0.55!(shield.west) + (0,-0.55)$)
        -| ([yshift=-8pt]shield.west);
    \node[edgelabel] at ($(state.east)!0.55!(shield.west) + (0,-0.72)$) {\(s_k\)};

    % --- Training path ---
    \draw[train] (state.south) -- ++(0,-0.35) -| (critic.north);
    \draw[train] (buffer.south) -- ++(0,-1.05) -| ([xshift=-6pt]critic.south);
    \node[edgelabel] at ($(buffer.south)!0.5!(critic.south) + (0.9,-0.55)$) {\(r_k,\phi,\psi\)};

    % --- Download feedback ---
    \draw[infer, gray!60] (client.south) -- ++(0,-0.45) -| ([xshift=5pt]buffer.south);

    % --- Group boxes ---
    \begin{pgfonlayer}{background}
        \node[groupbox, fill=gray!3, fit=(network)(manifest)(buffer),
              label={[font=\small\bfseries, text=gray!65]above:Environment}] {};
        \node[groupbox, fill=blue!2, fit=(actor)(critic), inner ysep=12pt,
              label={[font=\small\bfseries, text=blue!55!black]above:Lagrangian PPO}] {};
        \node[groupbox, fill=green!3, fit=(shield),
              label={[font=\small\bfseries, text=green!45!black]above:Runtime Shield}] {};
    \end{pgfonlayer}

    % --- Legend ---
    \node[font=\footnotesize, anchor=north] at (6.2, -2.35) {
        \tikz[baseline=-0.5ex]\draw[inferbold] (0,0) -- (0.55,0);~Inference \quad
        \tikz[baseline=-0.5ex]\draw[train] (0,0) -- (0.55,0);~Training only
    };

\end{tikzpicture}
}%
    \caption{Control loop. The base policy proposes \(a^{\mathrm{raw}}_k\); CPS banks on the VMAF knee and projects under a conformal throughput lower bound to produce \(a^{\mathrm{exec}}_k\). Training-only dashed paths (critic, dual variables) are optional for RL hosts.}
    \label{fig:arch}
\end{figure}

Decisions advance on the chunk clock, one step per playback duration \(T_{\mathrm{c}}\) (\(4\,\mathrm{s}\) in our harness)~\cite{mao2017neural,yin2015control}.

The two layers serve distinct roles. Optional learned-host training (Section~\ref{sec:system:cmdp}) handles stall and smoothness budgets in expectation over rollouts~\cite{paternain2019dual,achiam2017constrained}. However, it does not guarantee per-step feasibility under conformal throughput bounds. CPS provides this guarantee through decision-time banking and certified repair~\cite{alshiekh2018shielding,vovk2005algorithmic}.

\subsection{State, action, reward, and nonlinear perceptual mapping}

\emph{State \(s_k\).}
The observation holds buffer occupancy, past throughput samples, and the previous decision.
It also includes manifest lookahead: upcoming chunk sizes across the representation ladder, plus spatial complexity (SI) and temporal complexity (TI) indices~\cite{yin2015control,sodagar2011mpeg,mao2017neural}.
The VMAF ladder \(\{V_k(j)\}_{j\in\mathcal{B}}\) completes the observation, where \(V_k(j)\) is the perceptual score for chunk \(k\) at representation \(j\)~\cite{li2016vmaf}.
This last component is essential: without \(\{V_k(j)\}\), neither a learned host nor CPS can detect the content-dependent ladder saturation that bandwidth banking exploits.

\emph{Nonlinear, content-dependent fidelity.}
Saturation arises because quality does not scale linearly with rate: for each chunk \(k\), the map \(j\mapsto V_k(j)\) is not affine in the discrete index \(j\).
This map also varies across chunks. Different group-of-pictures (GOP) and texture structures yield different marginal VMAF gains between adjacent rungs. Therefore, the bitrate-to-VMAF relationship is both nonlinear and content-dependent.
Consequently, adjacent rungs may differ by a few VMAF points or by more than ten. Quality often saturates at the upper rungs of the ladder---precisely the structure that CPS banking exploits~\cite{li2016vmaf}. Perceptually aware, JND-based per-title encoding aims to remove this saturation at authoring time~\cite{menon2022perceptual}. In contrast, CPS exploits any saturation that remains in a deployed ladder at decision time.

\emph{Action.}
\(a^{\mathrm{raw}}_k\in\mathcal{B}\) is the policy's proposed representation index for chunk \(k\); the environment executes only \(a^{\mathrm{exec}}_k\).

\emph{QoE surrogate.}
Keeping to RobustMPC/Pensieve reporting conventions~\cite{yin2015control,mao2017neural}, we logged a scalar per-chunk surrogate with fidelity \(q_k \equiv V_k\!\left(a^{\mathrm{exec}}_k\right)\):
\begin{equation}
\label{eq:qoe-rep}
    \mathrm{QoE}^{\mathrm{(rep)}}_{k} = q_k - \mu\, \text{smooth}_k - \beta\, \tau^{\mathrm{stall}}_k,
\end{equation}
where \(\text{smooth}_k=\lvert q_k-q_{k-1}\rvert\) is the absolute VMAF change between adjacent chunks and \(\tau^{\mathrm{stall}}_k\) is the stall duration at step \(k\).
The first chunk takes a zero smoothness term.
Our shared harness fixed \(\beta{=}4.3\) per stall second and \(\mu{=}1.0\) per absolute VMAF-point change. The symbol \(\alpha\) remained reserved for the conformal miscoverage throughout.
The rebuffering weight is a Pensieve-style penalty, rescaled onto the VMAF fidelity axis.
Session QoE sums over chunks: \(\sum_{k\in\mathcal{K}_e}\mathrm{QoE}^{\mathrm{(rep)}}_{k}\) (Section~\ref{sec:eval}).
For a learned host, optional RL training substitutes a Lagrangian CMDP objective over rollouts (\ref{app:cmdp})~\cite{altman1999constrained,paternain2019dual}.

\subsection{Ladder saturation}
\label{sec:system:ladder}

The evaluated system reported one VMAF value per representation and video (pooled per-frame scores), so shield decisions operated on a session-mean ladder that did not vary with chunk index \(k\).
\tabref{tab:ladder:spacing} summarizes session-mean ladders for \(\LadNVideos\) publicly available reference titles (six Blender animations, six Xiph DERF clips); \figref{fig:ladder} illustrates per-chunk spread on a legacy four-title subset where frame-level VMAF logs were retained.
Across the twelve-title suite, \(\LadPctNonMonotone\%\) exhibited a non-monotone session-mean ladder. The mean top-rung gain (6000 vs.\ 2850\,kbps) was only \(\LadTopGainMean\) VMAF points. This pattern reflected strong saturation on animation titles (e.g.\ Big Buck Bunny at \(+0.03\)) but steeper climbs on natural content (e.g.\ Park Joy at \(+14.3\)).
Two properties made banking worthwhile. First, many adjacent pairs differed by at most one VMAF point on saturated titles, so a single-step downgrade cost almost nothing perceptually. Second, the top rungs saturated frequently, allowing the knee rule~\eqref{eq:knee} to bank bytes with no material quality loss.
For the monotone per-video ladder evaluated here, ranking feasible indices by VMAF produced the same choice as projection to the highest feasible index. CPS instead selects the smallest within-budget rung. This choice took effect only when saturation rendered downshifts near-lossless.
Chunk-level inversions appeared in \figref{fig:ladder} for the logged subset. However, banking did not depend on them. Section~\ref{sec:eval:perchunk} showed that the knee rule remained effective. In fact, it banked more when the shield observed the raw per-chunk ladder.

\begin{figure}[!htbp]
    \centering
    \evalfiglarge{figures/fig_ladder_perchunk.pdf}
    \caption{Representation ladder measured from the per-frame VMAF logs. (a)~Per-video pooled ladder (markers) against the range spanned by individual chunks (shaded). (b)~Distribution of the VMAF gap between adjacent representations across all \(\LadNChunks\) chunks; the dashed line marks a zero gap, below which the higher representation scores worse.}
    \label{fig:ladder}
\end{figure}

\begin{cps@table}
    \caption{Session-mean representation ladder across twelve evaluation titles.}
    \label{tab:ladder:spacing}
    \centering
    \evaltable{% Session-mean ladder summary (v19, twelve titles) — make_cps_paper_assets.py
\begin{tabularx}{\linewidth}{@{}>{\raggedright\arraybackslash}X c c c c c c@{}}
\toprule
\textbf{Video} & \textbf{Held-out} & \textbf{SI} & \textbf{TI} &
\textbf{Top gain} & \textbf{Span} & \textbf{Mono.} \\
\midrule
Big Buck Bunny & --- & 60.6 & 47.1 & +0.03 & 30.6 & Yes \\
Tears of Steel & --- & 41.0 & 31.8 & +4.09 & 26.1 & No \\
Sintel & \checkmark & 50.4 & 14.7 & +1.30 & 22.9 & Yes \\
Elephants Dream & --- & 51.8 & 27.0 & +7.87 & 31.2 & No \\
Ducks Take Off & --- & 94.9 & 16.6 & +11.59 & 36.6 & No \\
Cosmos Laundromat & \checkmark & 42.6 & 38.3 & +5.23 & 34.6 & No \\
Park Joy & --- & 119.7 & 34.8 & +14.31 & 46.5 & No \\
Old Town Cross & --- & 61.2 & 12.6 & +6.85 & 49.5 & Yes \\
Rush Hour & --- & 27.7 & 11.7 & +4.00 & 29.8 & No \\
Into Tree & --- & 47.4 & 14.4 & +9.12 & 36.1 & No \\
Sunflower & --- & 38.0 & 16.5 & +1.42 & 24.9 & No \\
Kristen \& Sara & \checkmark & 85.8 & 7.0 & +1.53 & 37.6 & Yes \\
\bottomrule
\end{tabularx}
}
    \tablenote{Top gain is VMAF(6000)\( - \)VMAF(2850) in kbps; Span is VMAF(6000)\( - \)VMAF(300). Mono.\ marks a non-decreasing session-mean ladder. Held-out titles (\checkmark) are excluded from PPO training. SI/TI are spatial/temporal complexity indices.}
\end{cps@table}

\subsection{Chunk-timescale dynamics}

\subsubsection{Throughput and latency}
Let \(\hat{C}_k>0\) be the trace-derived effective throughput, in kb/s, for chunk \(k\).
Every method saw the identical throughput trace in paired evaluations.
The choice of \(a^{\mathrm{exec}}_k\) sets the manifest payload size \(\sigma_{k}(a^{\mathrm{exec}}_k)\) in bits.
Write \(\tilde{C}_k=\max\{\hat{C}_k,\varepsilon_0\}\) for a small constant \(\varepsilon_0{>}0\) (the same numerical guard as in~\eqref{eq:dl}). The nominal download duration is \(d_k=\sigma_{k}(a^{\mathrm{exec}}_k)/(10^3\tilde{C}_k)\) seconds~\cite{yin2015control}.

\subsubsection{Buffer and stall}
Let \(B_k\ge 0\) be the buffered playback time when the download of chunk \(k\) begins~\cite{yin2015control,mao2017neural}.
Rebuffering (stall) and the buffer update follow
\begin{equation}
    \tau^{\mathrm{stall}}_k=\max\{0,\, d_k-B_k\},
\end{equation}
\begin{equation}
    \label{eq:buffer}
    B_{k+1}=\max\{0,B_k-d_k\}+T_{\mathrm{c}}.
\end{equation}
Randomness across episodes has a single source: trace and initial-condition sampling~\cite{fcc2015mba,kvalbein2014,bentaleb2018survey}.
Given \((s_k,a^{\mathrm{exec}}_k,\hat{C}_k)\), each transition is deterministic.

\subsection{Optional Lagrangian training for RL hosts}
\label{sec:system:cmdp}

When the base policy is learned, we trained PPO with a Lagrangian-relaxed constrained MDP (CMDP) objective that penalized rebuffering and smoothness alongside fidelity~\cite{altman1999constrained,paternain2019dual,achiam2017constrained}.
The dual variables adapt as training proceeds; \ref{app:cmdp} details the optional buffer-deviation and Lyapunov stabilizers~\cite{neely2010stochastic}.
None of this certifies per-chunk feasibility, which finite-sample PPO with clipped objectives cannot provide~\cite{schulman2017proximal}; CPS supplies this layer at deployment time.
Co-design wrapped the training environment in CPS, so the actor--critic observed the projected action distribution, whereas shield-at-eval-only policies trained without it.

% =========================================
\section{Method: Certified Perceptual Shield}
\label{sec:method}

CPS is a deterministic, model-agnostic wrapper around any base policy that proposes \(a^{\mathrm{raw}}_k\in\mathcal{B}\).
The rest of this section covers the conformal estimator, the banking and feasibility projection (Algorithm~\ref{alg:cps}), a predictive extension, optional PPO co-training, and the evaluation arms.

\subsection{Conformal throughput lower bound}
\label{sec:method:conformal}

The point predictor is the harmonic mean of the last \(\kappa\) observed throughputs, RobustMPC-style~\cite{yin2015control}.
After each step, the multiplicative residual \(r=C^{\mathrm{act}}/\hat{C}\) is appended to a rolling window of size \(W\).
Given \(n\) residuals, the split-conformal \(\alpha\)-quantile uses the finite-sample rank \(\lfloor\alpha(n{+}1)\rfloor\) (1-based). A linearly interpolated empirical quantile is excluded: it under-covers on short episodes~\cite{vovk2005algorithmic,angelopoulos2023gentle}.
If that rank is undefined---\(n\) too small---we fell back to a fixed scale \(\rho_{\mathrm{fb}}{=}0.80\).
The lower bound reads \(\underline{C}_k=\hat{C}_k\cdot q_{\alpha}\), clipped to stay at or below \(\hat{C}_k\).
Under residual exchangeability, the construction guarantees \(\mathbb{P}(C^{\mathrm{act}}\ge\underline{C}_k)\ge 1-\alpha\).
Given \(\underline{C}_k\), the certified download duration of rung \(j\) is
\begin{equation}
\label{eq:dl}
    d_k(j)=\min\!\left\{\frac{\sigma_k(j)}{\underline{C}_k\cdot 10^{3}+\varepsilon_0},\ 60\right\}.
\end{equation}

\begin{proposition}[Per-chunk stall certificate]
\label{prop:feas}
Fix chunk \(k\), executed index \(j\), buffer \(B_k\), and conformal lower bound \(\underline{C}_k\) computed with miscoverage \(\alpha\).
If \(d_k(j)\le \max\{B_k-m,0.1\}\) via~\eqref{eq:dl} and the realized throughput satisfies \(C^{\mathrm{act}}_k\ge \underline{C}_k\), then \(\tau^{\mathrm{stall}}_k=0\).
Under exchangeability of conformal residuals, \(\mathbb{P}(C^{\mathrm{act}}_k\ge \underline{C}_k)\ge 1-\alpha\) at each step where the bound is defined.
\end{proposition}
\begin{proof}[Proof sketch]
When \(C^{\mathrm{act}}_k\ge \underline{C}_k\), the realized download time is at most \(d_k(j)\), which by assumption does not exceed \(B_k\); hence~\eqref{eq:buffer} gives \(\tau^{\mathrm{stall}}_k=0\).
Coverage follows from standard split-conformal prediction with the \((n{+}1)\) rank correction~\cite{vovk2005algorithmic,angelopoulos2023gentle}.
\end{proof}

\emph{Scope of the guarantee.} The certificate provides marginal coverage over the exchangeable residual window: \(\mathbb{P}(C^{\mathrm{act}}_k\ge\underline{C}_k)\ge 1-\alpha\), averaged over draws. It does not provide coverage conditional on the current buffer or throughput regime. Distribution-free split conformal cannot do so without stronger assumptions~\cite{angelopoulos2023gentle}.

Abrupt \mbox{5G/mmWave} LOS/NLOS and scheduling shifts can break residual exchangeability~\cite{narayanan2021variegated}. The online defense is a sliding residual window of length \(W\): recent ratios both calibrate \(q_{\alpha}\) and immediately precede the predicted chunk. The bound therefore tracks local nonstationarity, even though finite-sample exactness holds only on an exchangeable window. If drift is large, coverage decays. The decay is bounded by a total-variation term between calibration and test residuals~\cite{barber2023beyond}. A weighted-conformal variant that down-weights stale residuals is a natural extension, which we leave to future work.
Empirically, we treated reported coverage as the operational check of Proposition~\ref{prop:feas}. On the primary synthetic \mbox{5G} pool, the certified arm met \(1-\alpha\) (\(\CPSGreedyCov\) vs.\ \(\CPScovTarget\)). Real broadband traces and tighter \(\alpha\) could sit marginally below the target (Section~\ref{sec:eval:improved}). This pattern matched the finite-sample and nonstationarity stress that the theory predicts.

\subsection{Banking and certified feasibility}
\label{sec:method:shield}

\emph{Safety set.} Indices that never exceed the proposal and fit the buffer margin \(m\) form
\begin{equation}
\label{eq:asafe}
    \mathcal{A}_{\mathrm{safe},k}(a)=\bigl\{j\le a:\ d_k(j)\le \max\{B_k-m,0.1\}\bigr\}.
\end{equation}
\emph{VMAF-knee banking.} Even when \(a\) is feasible, CPS selects the smallest rung within a perceptual budget \(\varepsilon\) of the proposal, namely
\begin{equation}
\label{eq:knee}
    j^{\mathrm{knee}}_k(a)=\min\bigl\{j\le a:\ V_k(a)-V_k(j)\le\varepsilon\bigr\},
\end{equation}
The shield then projects \(j^{\mathrm{knee}}_k(a)\) onto \(\mathcal{A}_{\mathrm{safe},k}(a)\); safety always takes precedence over \(\varepsilon\).
\emph{Remark.} Banking may run a lower rung than the proposal, and a lower rung can only shorten \(d_k(\cdot)\). Once~\eqref{eq:asafe} holds, safety projection therefore keeps the certificate intact.

Banking distinguishes CPS from prior VMAF-ranking shields, and the distinction is substantive. On the monotone ladders seen in practice, a ranking shield provably banks nothing.

\begin{proposition}[Ranking shields collapse on monotone ladders]
\label{prop:collapse}
Fix chunk \(k\) and suppose the VMAF ladder is nondecreasing in the rung index, \(V_k(0)\le\cdots\le V_k(L-1)\), and that download time \(d_k(\cdot)\) is nondecreasing in the rung index. Let \(\mathcal{F}=\{j: d_k(j)\le B_k-m\}\) be the feasible set. Then any shield that returns \(\arg\max_{j\in\mathcal{F}} V_k(j)\) returns the highest feasible index \(\max\mathcal{F}\), independently of any perceptual-loss budget \(\varepsilon\).
\end{proposition}
\begin{proof}
Because \(d_k\) is nondecreasing, \(\mathcal{F}\) is a prefix \(\{0,\dots,j_{\max}\}\) with \(j_{\max}=\max\mathcal{F}\). Since \(V_k\) is nondecreasing, \(\arg\max_{j\in\mathcal{F}} V_k(j)=j_{\max}\); the budget \(\varepsilon\) never enters the maximization.
\end{proof}
Ranking by retained VMAF is therefore equivalent to projection onto the highest feasible index, and it frees no bandwidth. CPS instead uses the proposal's quality and steps down to the smallest rung within \(\varepsilon\) of it~\eqref{eq:knee}. This mechanism lets CPS exploit rate--quality saturation: when top rungs carry near-identical VMAF, the knee lies well below \(j_{\max}\), and the freed bytes bank as buffer.

Algorithm~\ref{alg:cps} realizes~\eqref{eq:knee}--\eqref{eq:asafe}.
If \(B_k\le B_{\mathrm{crit}}\), CPS forces the lowest rung. Otherwise the algorithm computes the knee index via~\eqref{eq:knee} under the effective budget \(\varepsilon_{\mathrm{eff}}\) (equal to \(\varepsilon\) unless predictive mode is on; Section~\ref{sec:method:predictive}), then decrements the index until~\eqref{eq:dl} respects the margin \(m\).
The shield never upgrades a proposal.
Diagnostics log banked bits, VMAF surrendered, and whether realized throughput met \(\underline{C}_k\) (empirical coverage).

\begin{algorithm}[!t]
\footnotesize
\caption{Certified Perceptual Shield (compact)}
\label{alg:cps}
\begin{algorithmic}[1]
\STATE \textbf{Input:} \(a^{\mathrm{raw}}\), \(B_k\), \(\{\sigma_k(j),V_k(j)\}\), conformal estimator; parameters \((\varepsilon,m,B_{\mathrm{crit}})\)
\STATE \textbf{Output:} \((a^{\mathrm{exec}},\iota)\) with \(\iota=\mathbf{1}[a^{\mathrm{exec}}\neq a^{\mathrm{raw}}]\)
\STATE \(a\leftarrow\mathrm{clip}(a^{\mathrm{raw}})\); \(\underline{C}_k\leftarrow\) conformal lower bound; define \(d_k(\cdot)\) via~\eqref{eq:dl}
\IF{\(B_k\le B_{\mathrm{crit}}\)} \STATE \textbf{return} \((0,\mathbf{1}[a\neq0])\) \ENDIF
\STATE compute \(\varepsilon_{\mathrm{eff}}\) (Section~\ref{sec:method:predictive}); \(j\leftarrow j^{\mathrm{knee}}_k(a)\) via~\eqref{eq:knee} with \(\varepsilon_{\mathrm{eff}}\)
\WHILE{\(j>0\) and \(d_k(j)>\max\{B_k-m,0.1\}\)} \STATE \(j\leftarrow j-1\) \ENDWHILE
\STATE \(a^{\mathrm{exec}}\leftarrow\min\{a,j\}\); \textbf{return} \((a^{\mathrm{exec}},\mathbf{1}[a^{\mathrm{exec}}\neq a])\)
\end{algorithmic}
\end{algorithm}

\paragraph{Evaluation arms.}
\textsc{Raw}: we ran \(a^{\mathrm{raw}}\) unshielded.
\textsc{Safety}: Algorithm~\ref{alg:cps} with banking switched off (\(j^{\mathrm{knee}}{=}a\)), leaving conformal feasibility alone.
\textsc{Certified}: the full CPS, banking plus feasibility.
Paired seeds kept the Wilcoxon and TOST tests valid across arms.

\subsection{Predictive and risk-aware perceptual budget}
\label{sec:method:predictive}

The predictive budget widens when slack is scarce: it banks lightly when buffer headroom is ample and more aggressively when risk arises. While the buffer stays comfortable, \(\varepsilon_{\mathrm{eff}}{=}\varepsilon\) holds quality steady.
Two triggers widen the budget. The first is a drop of \(B_k\) below a risk threshold \(B_{\mathrm{risk}}\). The second is a short \(H\)-step look-ahead under a planning throughput floor (a low quantile of recent observations) that forecasts a dip below \(B_{\mathrm{risk}}\),
\begin{equation}
\label{eq:eps-eff}
    \varepsilon_{\mathrm{eff}}=\varepsilon+(\varepsilon_{\mathrm{risk}}-\varepsilon)\cdot f(B_k,\mathrm{lookahead}),
\end{equation}
with \(f\in[0,1]\). This lets the shield pre-bank more deeply before stalls force emergency cuts.
The ``improved'' experiments ran with \(\varepsilon{=}\CPSepsilon\), \(\varepsilon_{\mathrm{risk}}{=}4\), \(B_{\mathrm{risk}}{=}8\,\mathrm{s}\), and \(H{\ge}6\) at forecast quantile \(0.2\).

\subsection{Optional learned hosts and co-design}
\label{sec:method:codesign}

CPS does not require a learned policy.
When the host is an RL policy, we trained PPO (optionally Lagrangian-shaped) with Stable-Baselines~3~\cite{raffin2021sb3,schulman2017proximal} on content-aware observations: throughput history, buffer, SI/TI, VMAF ladder.
\emph{Co-design:} wrapping the training environment in CPS changed the action distribution observed by the actor--critic. The policy could then learn to propose bolder rungs that banking would trim for safety.
To gauge the effect, we compared two policies against each other under identical \mbox{5G} train/test provenance: one trained without the shield (shield-at-eval-only), one co-trained with CPS. Both faced the improved shield at test time.

\subsection{Implementation notes}
\label{sec:method:impl}

At the primary \mbox{5G} operating point the defaults were \(\alpha{=}\CPSalpha\) (coverage target \(\CPScovTarget\)), window \(W{=}200\), \(\kappa{=}5\), margin \(m{=}0.5\,\mathrm{s}\), \(B_{\mathrm{crit}}{=}0.3\,\mathrm{s}\), and buffer cap \(12\,\mathrm{s}\).
\tabref{tab:cps:hypers} collected the shield and predictive-banking defaults from the primary experiments. Empirical coverage was the fraction of steps with \(C^{\mathrm{act}}\ge\underline{C}_k\).
Per decision, the projection ran in \(O(L{+}H)\) over the \(L\) ladder rungs and an \(H\)-step look-ahead. Conformal-quantile upkeep on the length-\(W\) residual window was \(O(1)\) amortized.
Across \(20{,}000\) rollout decisions on a commodity CPU (\(L{=}6\), full predictive shield), per-decision wall time was \(131\,\mu\mathrm{s}\) at the median and \(310\,\mu\mathrm{s}\) at the \(99\)th percentile. This is roughly \(10^{-5}\) of the \(4\,\mathrm{s}\) chunk budget. CPS therefore ran inline, and the control loop was unaffected.

\begin{cps@table}
    \caption{CPS default hyperparameters (primary \mbox{5G} suite).}
    \label{tab:cps:hypers}
    \centering
    \evaltable{%
    \begin{tabular}{@{}ll@{}}
    \toprule
    \textbf{Parameter} & \textbf{Value} \\
    \midrule
    Perceptual budget \(\varepsilon\) & \(\CPSepsilon\) VMAF point \\
    Conformal miscoverage \(\alpha\) & \(\CPSalpha\) (target coverage \(\CPScovTarget\)) \\
    Residual window \(W\) / predictor \(\kappa\) & 200 / 5 \\
    Buffer margin \(m\) / critical \(B_{\mathrm{crit}}\) & \(0.5\,\mathrm{s}\) / \(0.3\,\mathrm{s}\) \\
    Risk budget \(\varepsilon_{\mathrm{risk}}\) / \(B_{\mathrm{risk}}\) & \(4.0\) / \(8\,\mathrm{s}\) \\
    Look-ahead horizon \(H\) / forecast quantile & \(\ge 6\) / 0.2 \\
    \bottomrule
    \end{tabular}%
    }
\end{cps@table}

% =========================================
\section{Experimental evaluation}
\label{sec:eval}

We organized the evaluation around six research questions:

\begin{itemize}
  \item[\textbf{RQ1}] \emph{Banking vs.\ safety.} Does VMAF-knee banking reduce bitrate and rebuffering relative to conformal safety-only projection, with VMAF inside \(\pm\varepsilon\)?
  \item[\textbf{RQ2}] \emph{Coverage.} Does empirical conformal coverage meet the finite-sample target \(1{-}\alpha\)?
  \item[\textbf{RQ3}] \emph{Predictive banking.} Does a risk-aware, look-ahead budget deepen savings without leaving the perceptual envelope?
  \item[\textbf{RQ4}] \emph{Co-design.} How does training with CPS in the loop change the certified operating point relative to shield-at-eval-only operation?
  \item[\textbf{RQ5}] \emph{Per-chunk robustness.} Does banking remain sound and near-lossless when the shield operates on the true (non-monotone) per-chunk ladder rather than the pooled session-mean ladder?
  \item[\textbf{RQ6}] \emph{Hyperparameter sensitivity.} Are the results stable across the perceptual budget \(\varepsilon\) and miscoverage \(\alpha\), and does coverage track the target \(1{-}\alpha\)?
\end{itemize}

\subsection{Experimental setup and datasets}
\label{sec:eval:setup}

The testbed was a multi-video, chunk-timescale emulator (\(K_e{=}48\) chunks/episode, \(T_{\mathrm{c}}{=}4\,\mathrm{s}\), buffer cap \(12\,\mathrm{s}\) unless noted). It drove twelve publicly available reference titles with session-mean VMAF ladders~\cite{li2016vmaf}: six Blender open animations (Big Buck Bunny, Tears of Steel, Sintel, Elephants Dream, Ducks Take Off, Cosmos Laundromat) and six Xiph DERF clips (Park Joy, Old Town Cross, Rush Hour, Into Tree, Sunflower, Kristen \& Sara).
We selected them for ladder structure, not catalogue breadth: saturation-rich animation (Big Buck Bunny, Tears of Steel), steep natural motion (Park Joy, Ducks Take Off), wide-span DERF clips (Into Tree, Old Town Cross), and mixed held-out titles (Sintel, Cosmos Laundromat, Kristen \& Sara) for generalization checks (\tabref{tab:ladder:spacing}).
Nine titles appeared in PPO co-design training; three were held out entirely (Sintel, Cosmos Laundromat, Kristen \& Sara).
The negative broadband result (Section~\ref{sec:eval:improved}) showed that we did not select regimes that favor banking. CPS is a deterministic per-chunk projection fixed by the observed VMAF ladder and buffer state, so ladder shape---not semantic content---drove its effect. The twelve-title suite therefore spanned both saturated and steep/inverted session-mean ladders (\(\LadPctNonMonotone\%\) non-monotone).

\textbf{Primary \mbox{5G} suite.} The primary pool comprised synthetic \mbox{5G} traces with intermittent dips and outages, held out from the co-design training seeds (\(\CPSepisodes\) paired episodes). We generated them synthetically because public per-chunk \mbox{5G/mmWave} throughput logs suitable for chunk-level ABR replay remained scarce. For realism, the generator was tuned to reported \mbox{5G/mmWave} dip-and-outage dynamics~\cite{narayanan2021variegated,fezeu2024midband}, and the broadband suite below added a real-network counterpart.
Train and test trace pools stayed source-disjoint. Co-design policies trained on the training pool alone and faced held-out test seeds.

\textbf{Broadband suite.} The broadband suite used real-measured traces: the FCC Measuring-Broadband-America fixed-line logs~\cite{fcc2015mba} and the Norway/Nornet mobile-broadband mobility logs~\cite{kvalbein2014} (27 held-out test traces, source-disjoint from training). Both are widely used public datasets in ABR research. They provided a real-network counterpart to the synthetic \mbox{5G} suite and showed where banking offered limited headroom.

\textbf{Hosts.} Five host families covered the ABR design space: bitrate-greedy (always at the top rung), buffer-based ABR (BBA-like), BOLA (the Lyapunov buffer-based rule of the dash.js reference player~\cite{spiteri2018bola,spiteri2018dashjs}), robust model-predictive control (RobustMPC~\cite{yin2015control}), and PPO policies (content-aware ``proposed'' and content-blind Pensieve-style). RobustMPC was a five-segment planner over a throughput prediction discounted by recent worst-case error. CPS carried no host bias: one wrapper, applied at inference across all five families, tuned to none of them.

\textbf{Arms and pairing.} Every host rolled out \textsc{Raw}, \textsc{Safety}, and \textsc{Certified} on identical Python seeds---the same video--trace pairs.
Pairing fixes the random seed and the \((\text{video},\text{trace})\) assignment across arms, so metrics differ only because of shielding, not because of resampling.
Each paired contrast is the per-episode metric difference on the same \((\text{video},\text{trace})\) seed across arms (e.g., \textsc{Certified} minus \textsc{Safety} on episode \(i\)).
We reported mean delivered bitrate, total rebuffering, stall-free rate, mean VMAF, intervention rate, and empirical conformal coverage. Significance rested on paired Wilcoxon signed-rank tests (zero-method \texttt{zsplit}) and TOST equivalence of VMAF within \(\pm\varepsilon{=}\CPSepsilon\)~\cite{lakens2017equivalence}. Defaults were \(\alpha{=}\CPSalpha\) and coverage target \(\CPScovTarget\).
We set \(\varepsilon{=}\CPSepsilon\) as both the knee budget and the TOST margin. One VMAF point is a small fraction of typical adjacent-rung and top-rung gaps on our ladders (\tabref{tab:ladder:spacing}) and is conservative relative to perceptually aware per-title encoding practice~\cite{menon2022perceptual}. It is an engineering design choice aligned with~\eqref{eq:knee}, not a preregistered clinical equivalence margin~\cite{lakens2017equivalence}.
The \(p\)-values we quoted were raw and two-sided. Since several paired contrasts ran at once (bitrate and rebuffering across the host families and two shield settings), we checked that every reported bitrate and rebuffering effect survived Holm--Bonferroni correction at \(0.05\) over the reported family. The paired-episode bootstrap additionally supported the co-design VMAF and rebuffering contrasts.
A caveat on diversity followed. Each episode was one \((\text{video},\text{trace})\) pair---video drawn uniformly from the twelve titles, trace from the held-out pool. Observations therefore clustered by title, and effective independent-content diversity was twelve, not \(\CPSepisodes\).

We treated the paired tests as evidence of within-pair shield effects---the shielded arm vs.\ the unshielded arm on the same video--trace pair---and we reported per-title ladder structure (\tabref{tab:ladder:spacing}) with held-out breakdown where applicable. Vanishingly small \(p\)-values (say \(<10^{-4}\)) tracked sign-consistency of a paired difference over \(\CPSepisodes\) pairs. They constituted strong evidence for a within-pair effect, not precise magnitudes.

\subsection{Model-agnostic banking on \mbox{5G} (RQ1--RQ2)}
\label{sec:eval:main}

\tabref{tab:cps:full} reported absolute metrics for all three arms---\textsc{Raw}, \textsc{Safety}, and \textsc{Certified}---on synthetic \mbox{5G} test traces for five host families: greedy, BBA, BOLA, RobustMPC, and a Pensieve-style negative control.
\figref{fig:cps:overview} distilled only the certified-vs.-safety banking contrast in bar form (not \textsc{Raw}); the table lists the underlying episode means and significance tests for every arm.
Banking behaved consistently across all four deterministic hosts on the twelve-title suite. Certified banking vs.\ safety saved 4.5\%, 4.1\%, 4.1\%, and 3.8\% bitrate on greedy, BBA, BOLA, and RobustMPC, respectively, and reduced rebuffering by 9.3\%, 8.9\%, 8.9\%, and 9.2\% (\tabref{tab:cps:full}; greedy \(p{=}\CPSGreedyPRebCs\)).
Steep and non-saturated ladders (\tabref{tab:ladder:spacing}) kept absolute savings modest---the saturation-dependent regime the mechanism predicts---yet every host stayed perceptually near-lossless.
Partitioning the greedy certified--safety contrast by session-mean top-rung VMAF gain (\tabref{tab:ladder:spacing}) made that dependence explicit. On titles with top gain at most \(2\) VMAF points (Big Buck Bunny, Sintel, Sunflower, Kristen \& Sara; \(n{=}60\) episodes) banking saved \(4.6\%\) bitrate and \(9.9\%\) rebuffering. On steep titles with top gain at least \(9\) (Park Joy, Ducks Take Off, Into Tree; \(n{=}51\)) the same contrast shrank to \(2.1\%\) bitrate and \(3.7\%\) rebuffering. The pooled greedy headline (4.5\% bitrate saved and 9.3\% rebuffering reduced) therefore mixed both regimes rather than reflecting a saturation-only subset.
Mean VMAF differences were perceptually negligible, at \(\CPSGreedyVMcs\), \(\CPSBbaVMcs\), \(\CPSBolaVMcs\), and \(\CPSMpcVMcs\) points across greedy, BBA, BOLA, and RobustMPC, with every value inside \(\pm\varepsilon\) by TOST.
The two shield layers split the work. Against \textsc{Raw}, certified greedy reduced rebuffering by 55.5\% (BBA: 51.0\%). Conformal safety eliminated the catastrophic stalls; banking added the rest of the headroom.
\figref{fig:cps:tradeoff} and \figref{fig:cps:cdf} showed per-episode operating points and rebuffering CDFs for greedy and BBA.
Coverage followed a distinct pattern on the synthetic \mbox{5G} pool. Empirical coverage there was \(\CPSGreedyCov\) (greedy) and \(\CPSBbaCov\) (BBA) against a \(\CPScovTarget\) target, at or above the nominal level on synthetic \mbox{5G}, consistent with the \((n{+}1)\) conformal calibration at these episode lengths (Section~\ref{sec:method:conformal}).
Because coverage depends on the throughput trace and the conformal calibrator rather than the base policy, it was identical (\(\CPSGreedyCov\)) across every host and shield variant on the \mbox{5G} test set (\figref{fig:cps:coverage}).
Proposition~\ref{prop:feas} is a per-chunk probabilistic certificate, not a per-episode stall-free one. At coverage \(\CPSGreedyCov\), about \(8\%\) of chunks could have realized throughput fall below \(\underline{C}_k\). The forced downshift at \(B_{\mathrm{crit}}\) and the first-chunk cold start explained the residual rebuffering in the certified arm.
Pensieve, the negative control, degraded to near-floor bitrates with near-zero stalls, so CPS rarely intervened. Banking could not rescue a policy that never requested saturated rungs.

That high intervention rate under greedy (\tabref{tab:cps:full}) was fixed by construction; it did not imply that the shield dominated an informed base policy. Greedy proposed the top rung every time, so nearly every banking downshift or feasibility repair appeared as an intervention. The rate gauged how far a host sat from the certified operating point, not how much of the outcome to attribute to the shield. This is why we ran the model-agnostic contrast across hosts spanning the spectrum.

Certified intervention decreased across all hosts: \(\CPSGreedyCertInt\%\) for always-top greedy, \(\CPSBolaCertInt\%\) for BOLA, \(\CPSMpcCertInt\%\) for RobustMPC---which already planned conservatively against a discounted throughput estimate---and near zero for the conservative Pensieve-style control. The sharper the base policy, the lighter the shield's load. Coverage nevertheless held at \(\CPSGreedyCov\) throughout. The certified metrics therefore depended increasingly on the base policy's own decisions, with intervention only where banking headroom remained.

\begin{figure}[!htbp]
    \centering
    \evalfig{figures/fig_cps_overview.pdf}
    \caption{CPS banking overview (certified vs.\ safety, \(\CPSepisodes\) paired episodes, \(\varepsilon{=}\CPSepsilon\)). For two synthetic \mbox{5G} hosts and a broadband host, bars show percent bitrate saved, percent rebuffering reduced, and mean VMAF cost (mean VMAF difference in points, \textsc{Certified} minus \textsc{Safety}). Banking is significant and perceptually near-lossless on \mbox{5G} (magnitude scales with ladder saturation) but collapses on broadband, where top rungs are rarely jointly feasible and saturated.}
    \label{fig:cps:overview}
\end{figure}

\begin{figure}[!htbp]
    \centering
    \evalfig{figures/fig_cps_coverage.pdf}
    \caption{Per-episode conformal coverage of the certified arm on the synthetic \mbox{5G} test pool (\(\alpha{=}\CPSalpha\), target \(\CPScovTarget\), dashed). Coverage depends only on the throughput trace and calibrator, so the distribution is host-invariant by construction on this pool (RQ2).}
    \label{fig:cps:coverage}
\end{figure}

\begin{figure}[!htbp]
    \centering
    \evalfig{figures/fig_cps_tradeoff.pdf}
    \caption{Per-episode operating points on synthetic \mbox{5G} (paired seeds). Diamonds mark arm means. Safety removes catastrophic stalls relative to \textsc{Raw}; certified banking further reduces rebuffering at nearly unchanged mean VMAF.}
    \label{fig:cps:tradeoff}
\end{figure}

\begin{figure}[!htbp]
    \centering
    \evalfig{figures/fig_cps_cdf_reb.pdf}
    \caption{Empirical CDFs of episode rebuffering for greedy and BBA hosts. The certified arm stochastically dominates safety and raw on the stall axis.}
    \label{fig:cps:cdf}
\end{figure}

\begin{cps@table}
    \caption{Paired arms on synthetic \mbox{5G} test traces (\(\CPSepisodes\) episodes, \(\varepsilon{=}\CPSepsilon\), \(\alpha{=}\CPSalpha\)). Reb.\ is total rebuffering; Cover.\ is empirical conformal coverage; Interv.\ is shield intervention rate.}
    \label{tab:cps:full}
    \centering
    \evaltable{% Full paired arms on 5G test (macros from make_cps_paper_assets.py).
\begin{tabularx}{\linewidth}{@{}>{\raggedright\arraybackslash}X c ccccc@{}}
\toprule
\textbf{Host} & \textbf{Arm} & \textbf{Reb.} & \textbf{VMAF} & \textbf{Bitrate} & \textbf{Cover.} & \textbf{Interv.} \\
\midrule
Greedy & \textsc{Raw} & \CPSGreedyRawReb{}\,s & \CPSGreedyRawVM{} & \CPSGreedyRawBr{} & --- & 0.0\\
 & \textsc{Safety} & \CPSGreedySafReb{}\,s & \CPSGreedySafVM{} & \CPSGreedySafBr{} & \CPSGreedySafCov{} & \CPSGreedySafInt{}\%\\
 & \textsc{Certified} & \CPSGreedyCertReb{}\,s & \CPSGreedyCertVM{} & \CPSGreedyCertBr{} & \CPSGreedyCertCov{} & \CPSGreedyCertInt{}\%\\
\midrule
BBA & \textsc{Raw} & \CPSBbaRawReb{}\,s & \CPSBbaRawVM{} & \CPSBbaRawBr{} & --- & 0.0\\
 & \textsc{Safety} & \CPSBbaSafReb{}\,s & \CPSBbaSafVM{} & \CPSBbaSafBr{} & \CPSBbaSafCov{} & \CPSBbaSafInt{}\%\\
 & \textsc{Certified} & \CPSBbaCertReb{}\,s & \CPSBbaCertVM{} & \CPSBbaCertBr{} & \CPSBbaCertCov{} & \CPSBbaCertInt{}\%\\
\midrule
BOLA & \textsc{Raw} & \CPSBolaRawReb{}\,s & \CPSBolaRawVM{} & \CPSBolaRawBr{} & --- & 0.0\\
 & \textsc{Safety} & \CPSBolaSafReb{}\,s & \CPSBolaSafVM{} & \CPSBolaSafBr{} & \CPSBolaSafCov{} & \CPSBolaSafInt{}\%\\
 & \textsc{Certified} & \CPSBolaCertReb{}\,s & \CPSBolaCertVM{} & \CPSBolaCertBr{} & \CPSBolaCertCov{} & \CPSBolaCertInt{}\%\\
\midrule
RobustMPC & \textsc{Raw} & \CPSMpcRawReb{}\,s & \CPSMpcRawVM{} & \CPSMpcRawBr{} & --- & 0.0\\
 & \textsc{Safety} & \CPSMpcSafReb{}\,s & \CPSMpcSafVM{} & \CPSMpcSafBr{} & \CPSMpcSafCov{} & \CPSMpcSafInt{}\%\\
 & \textsc{Certified} & \CPSMpcCertReb{}\,s & \CPSMpcCertVM{} & \CPSMpcCertBr{} & \CPSMpcCertCov{} & \CPSMpcCertInt{}\%\\
\midrule
Pensieve & \textsc{Raw} & \CPSPenRawReb{}\,s & \CPSPenRawVM{} & \CPSPenRawBr{} & --- & 0.0\\
 & \textsc{Safety} & \CPSPenSafReb{}\,s & \CPSPenSafVM{} & \CPSPenSafBr{} & \CPSPenSafCov{} & \CPSPenSafInt{}\%\\
 & \textsc{Certified} & \CPSPenCertReb{}\,s & \CPSPenCertVM{} & \CPSPenCertBr{} & \CPSPenCertCov{} & \CPSPenCertInt{}\%\\
\bottomrule
\end{tabularx}
}
    \tablenote{Certified--safety Wilcoxon: greedy bitrate \(p{=}\CPSGreedyPBWcs\), rebuffering \(p{=}\CPSGreedyPRebCs\); BBA bitrate \(p{=}\CPSBbaPBWcs\), rebuffering \(p{=}\CPSBbaPRebCs\). Improved predictive banking defaults in \tabref{tab:cps:hypers}.}
\end{cps@table}

\subsection{Predictive banking and broadband boundary (RQ3)}
\label{sec:eval:improved}

With risk-aware \(\varepsilon(B)\) and dip forecasting enabled, certified banking vs.\ safety deepened savings to a 14.6\% bitrate reduction (greedy) and 12.5\% (BBA), and cut rebuffering by 12.6\% and 13.0\%, respectively.
By relaxing the perceptual budget ahead of forecast throughput dips, the shield can bank earlier while staying inside the \(\varepsilon\) envelope; mean VMAF changes remained negligible (\(\CPSGreedyImpVMcs\), \(\CPSBbaImpVMcs\) points).

The broadband suite marked the edge of these gains. On source-disjoint real-measured broadband traces (FCC MBA and Norway/Nornet~\cite{fcc2015mba,kvalbein2014}), greedy certified--safety bitrate savings shrank to 0.9\%. Top rungs were seldom jointly feasible and saturated, so banking had little remaining headroom.
Even so, safety vs.\ raw still reduced rebuffering by 98.7\%. Under broadband stress the conformal feasibility layer---not banking---accounted for most of the improvement, because saturated top rungs left almost no knee headroom for proactive banking.
On the same real-measured traces at \(\alpha{=}\CPSalpha\), empirical conformal coverage was \(\CPSbbCov\) (below the \(\CPScovTarget\) target at \(n{=}204\) episodes), whereas the synthetic \mbox{5G} pool reached \(\CPSGreedyCov\) (\tabref{tab:cps:full}, \figref{fig:cps:coverage}).
The ``BB greedy'' group in \figref{fig:cps:overview} made the regime split explicit: banking savings were largest on synthetic \mbox{5G}.
This split was by design: CPS banking paid off precisely when the ladder knee sat within reach, and broadband placed it out of reach.

\paragraph{Coverage under injected drift.}
To stress-test exchangeability without new traces, we replayed the synthetic \mbox{5G} pool with a mid-episode throughput shock (chunks \(28\)--\(37\): multiply trace throughput by \(0.3\), mimicking a sudden line-of-sight (LOS) loss or NLOS dip).
After a \(20\)-chunk conformal warm-up, per-chunk coverage in the steady pre-shock window was \(0.928\) (\(n{=}1{,}632\) chunks).
During the shock it was \(0.929\) (\(n{=}2{,}040\)); in the post-shock recovery window it rose to \(0.944\) (\(n{=}2{,}040\)), all against a \(1{-}\alpha{=}0.90\) target (greedy certified arm, \(\CPSepisodes\) episodes).
\figref{fig:cps:coverage_drift} plotted the rolling chunk coverage; the bound tracked the shock without collapsing below the target, consistent with the sliding-window defense in Section~\ref{sec:method:conformal}.

\begin{figure}[!htbp]
    \centering
    \evalfig{figures/fig_cps_coverage_drift.pdf}
    \caption{Rolling per-chunk conformal coverage under an injected throughput shock (shaded; \(0.3\times\) scale, chunks \(28\)--\(37\)) on the synthetic \mbox{5G} pool. Dashed line: target \(1{-}\alpha{=}\CPScovTarget\).}
    \label{fig:cps:coverage_drift}
\end{figure}

\subsection{Learned hosts and co-design (RQ4)}
\label{sec:eval:codesign}

Two PPO settings warranted separate reporting. The primary host was a content-aware policy trained on source-disjoint broadband traces and evaluated, without further training, on \mbox{5G}---the model-agnostic deployment case. The co-design pair comprised two policies trained on matched \mbox{5G} traces---one shield-at-eval-only, one co-trained with CPS in the loop---compared on their certified arms (\tabref{tab:cps:codesign} and \figref{fig:cps:codesign}).
On the primary broadband-trained policy, certified banking vs.\ safety saved 35.3\% bitrate and reduced rebuffering by 79.3\% (\(p{=}\CPSPropPRebCs\)), at VMAF \(\Delta{=}\CPSPropVMcs\) (TOST within \(\pm\varepsilon\)).
The effect exceeded that of greedy or BBA. The learned policy still proposed high rungs often on the twelve-title suite, which left more saturation headroom for CPS.
Under the improved shield, banking saved 21.9\% bitrate and reduced rebuffering by 29.6\%, mirroring the deeper predictive-banking gains on deterministic hosts in Section~\ref{sec:eval:improved}.

\emph{Co-design.} Both policies trained on matched \mbox{5G} traces (source-disjoint test). Under the improved shield at eval, the co-trained policy's certified arm attained VMAF \(\CPScoCpsVM\) and bitrate \(\CPScoCpsBr\)\,kb/s, against \(\CPScoRegVM\)/\(\CPScoRegBr\) for shield-at-eval-only operation.
The \(\Delta\) row in \tabref{tab:cps:codesign} was the difference of those certified-arm means (\(+\CPScoDVM\) VMAF, \(+\CPScoDReb\)\,s rebuffering), and it matched the raw rows above (\(\CPScoCpsVM-\CPScoRegVM\), \(\CPScoCpsReb-\CPScoRegReb\)).
The fidelity gain could be weighed against its stall cost through \(\mathrm{QoE}(w)=\overline{\mathrm{VMAF}}-w\cdot\mathrm{rebuffer(s)}\). The mean-QoE crossover was \(w^{\star}{=}\CPScoWstar\) VMAF points per second (ratio of the \(\Delta\) row entries): co-design favored fidelity-oriented weights \(w<w^{\star}\); shield-at-eval-only prevailed where stall aversion was stronger.
Intuitively, \(w^{\star}\) is the stall penalty at which the co-trained fidelity gain and extra rebuffering exactly balance in mean QoE; weights above \(w^{\star}\) favor the shield-at-eval-only policy.
Paired Wilcoxon tests and episode-bootstrap intervals were reported as \(\CPScoPVM\) (VMAF) and \(\CPScoPReb\) (rebuffering).
The bootstrap CI was \([\CPScoDVMlo,\CPScoDVMhi]\), with stall cost \([\CPScoDRebLo,\CPScoDRebHi]\)\,s and crossover \([\CPScoWstarLo,\CPScoWstarHi]\) (median \(\CPScoWstarMed\); positive VMAF in \(\CPScoDVMposPct\%\) of resamples).
Co-design remained a supporting result under a single training seed (Section~\ref{sec:limits}).

\begin{figure}[!htbp]
    \centering
    \evalfig{figures/fig_cps_codesign.pdf}
    \caption{Co-design on \mbox{5G} test (certified arm, improved shield). Left: per-episode VMAF--rebuffer points (diamond markers = episode means). Right: mean \(\Delta\mathrm{QoE}(w)\) vs.\ stall weight; the dashed line marks the crossover \(w^{\star}{=}\CPScoWstar\).}
    \label{fig:cps:codesign}
\end{figure}

\begin{cps@table}
    \caption{Co-design on \mbox{5G} test (certified arm, improved shield at eval).}
    \label{tab:cps:codesign}
    \centering
    \evaltable{% Co-design: certified arm only (5G test, improved shield at eval).
\begin{tabularx}{\linewidth}{@{}>{\raggedright\arraybackslash}X cccc@{}}
\toprule
\textbf{Policy} & \textbf{VMAF} & \textbf{Reb. (s)} & \textbf{Bitrate} & \textbf{Cover.} \\
\midrule
Eval-only shield & \CPScoRegVM & \CPScoRegReb & \CPScoRegBr & \CPScoRegCov \\
Co-trained (CPS) & \CPScoCpsVM & \CPScoCpsReb & \CPScoCpsBr & \CPScoCpsCov \\
\midrule
$\Delta$ (co $-$ eval) & $+\CPScoDVM$ & $+\CPScoDReb$ & --- & --- \\
\bottomrule
\end{tabularx}
}
    \tablenote{The \(\Delta\) row equals the certified-arm difference in the rows above. Paired Wilcoxon: VMAF \(p{=}\CPScoPVM\), rebuffering \(p{=}\CPScoPReb\); bootstrap \(\Delta\)VMAF CI \([\CPScoDVMlo,\CPScoDVMhi]\) over \(\CPScoBootN\) episodes (requires synced \texttt{episodes.csv}).}
\end{cps@table}

\subsection{Per-chunk (non-monotone) ladder ablation (RQ5)}
\label{sec:eval:perchunk}

The headline runs applied the knee rule to the pooled session-mean ladder, which is monotone. We therefore tested whether banking remained effective on the true per-chunk ladders, which were often non-monotone (\(\PCinvPct\%\) of the \(\PCnChunks\) chunks held at least one inverted adjacent pair; \tabref{tab:ladder:spacing}).
To address this, we replayed the deterministic banking stage on the per-chunk logs. We compared a shield that saw only the pooled ladder with one that saw the per-chunk ladder (\tabref{tab:perchunk}, \figref{fig:cps:perchunk}). Banking alone depended on VMAF; the feasibility projection ignored it and stayed identical in both cases, which isolated the effect of the pooling simplification.
Three findings held at the default \(\varepsilon{=}\PCeps\), spanning correctness, savings, and pooling error.
First, the knee rule remained well defined on inverted ladders and banked more, not less: mean saved bitrate was \(\PCbwPerchunk\%\) per-chunk against \(\PCbwPooled\%\) pooled, because a lower rung that inverted above the proposal satisfied the budget immediately---inversions widened banking opportunity rather than breaking the rule.
Second, a per-chunk shield held realized per-chunk perceptual cost inside budget by construction (mean \(\PCcostPerchunk\) VMAF points, never above \(\varepsilon\)), so the ablation validated chunk-wise enforcement when true ladders were available.
Third, pooling was not free. Judged against the true per-chunk VMAF, the pooled-ladder shield overran the budget on \(\PCviolPooled\%\) of chunks (mean cost \(\PCcostPooled\) points), which explains why session-mean evaluation could understate perceptual spend even when banking bytes looked similar.
Taken together, the pooled-ladder headline numbers were, if anything, conservative with respect to banking savings.
Per-chunk operation was sound: it eliminated residual budget violations at the cost of a per-chunk VMAF feed, and it aligned the content-diversity motivation with the pooled evaluation.

\begin{table}[t]
  \centering
  \caption{Banking-stage ablation on real per-chunk VMAF ladders (greedy proposal, \(\PCnChunks\) chunks). A per-chunk shield banks slightly more bytes than the pooled shield while keeping realized per-chunk cost \(\le\varepsilon\) by construction; the pooled shield incurs a small per-chunk budget-violation rate. Feasibility is VMAF-independent and unchanged.}
  \label{tab:perchunk}
  \small
  \evaltable{% Auto-generated by ablation_perchunk_ladder.py
\begin{tabular}{lcccccc}
\toprule
$\varepsilon$ & \multicolumn{2}{c}{Mean knee rung} & \multicolumn{2}{c}{Bitrate saved (\%)} & \shortstack{Per-chunk cost\\(pts)} & \shortstack{Budget viol.\\pooled (\%)} \\
\cmidrule(lr){2-3}\cmidrule(lr){4-5}
 & pooled & per-chunk & pooled & per-chunk & per-chunk & \\
\midrule
0.5 & 4.58 & 4.34 & 22.3 & 28.1 & 0.02 & 8.8 \\
1.0 & 4.15 & 3.94 & 44.6 & 45.8 & 0.32 & 11.5 \\
2.0 & 3.60 & 3.39 & 58.2 & 59.7 & 0.98 & 18.6 \\
4.0 & 2.93 & 2.88 & 68.7 & 68.7 & 1.89 & 7.1 \\
\bottomrule
\end{tabular}
}
\end{table}

\begin{figure}[t]
  \centering
  \evalfig{figures/fig_cps_perchunk.pdf}
  \caption{Per-chunk vs.\ pooled ladder banking across the perceptual budget \(\varepsilon\). Left: banked bitrate is comparable (slightly larger per-chunk). Right: realized per-chunk perceptual cost stays under \(\varepsilon\) for the per-chunk shield; the pooled shield tracks it closely on average but violates the budget on a minority of chunks.}
  \label{fig:cps:perchunk}
\end{figure}

\subsection{Sensitivity to \(\varepsilon\) and \(\alpha\) (RQ6)}
\label{sec:eval:sensitivity}

The headline runs fixed the pair at \(\varepsilon{=}\CPSepsilon\) and \(\alpha{=}\CPSalpha\). \tabref{tab:abl:eps}, \tabref{tab:abl:alpha}, and \figref{fig:cps:ablation} replayed the same greedy host on the synthetic \mbox{5G} test pool at \(\CPSepisodes\) paired episodes (twelve-title roster).
We first considered the perceptual budget. Sweeping \(\varepsilon\in\{\AblEpsLo,\dots,\AblEpsHi\}\) (at \(\alpha{=}\CPSalpha\)) traded banking for fidelity.
Bitrate saved vs.\ \textsc{Safety} climbed from 2.0\% to 24.6\%, while the mean VMAF change remained negligible (at most \(\AblVmMaxAbs\) points) and never left the matching \(\pm\varepsilon\) envelope---so wider budgets traded fidelity slack for banking headroom in a predictable way.
Empirical conformal coverage remained invariant with \(\varepsilon\). Coverage depended on the throughput trace and calibrator alone, not the banking budget, which showed that banking and the feasibility certificate operated independently by design.
We next considered miscoverage. Sweeping \(\alpha\in\{0.05,0.10,0.20\}\) (at \(\varepsilon{=}\CPSepsilon\)) on the same synthetic \mbox{5G} pool moved empirical coverage in step with the finite-sample target \(1{-}\alpha\): from \(\AblCovMax\) at target \(0.95\) through the default \(0.919\) at target \(0.90\) to \(\AblCovMin\) at target \(0.80\).
Tighten \(\alpha\) and both coverage and intervention rate rose; loosen it and both decreased, as the \((n{+}1)\) construction predicts.
At \(\alpha{=}0.05\) the mean was \(0.945\), marginally below the \(0.95\) target at \(n{=}204\) (Section~\ref{sec:eval:improved}). The default operating point therefore lay on a smooth, well-behaved frontier rather than a narrowly tuned optimum, with headline savings stable once \(\varepsilon\) and \(\alpha\) were fixed.

\begin{cps@table}
  \centering
  \caption{Sensitivity to the perceptual budget \(\varepsilon\) (greedy host, \mbox{5G} test, \(\alpha{=}\CPSalpha\), certified vs.\ safety). Banking scales smoothly with \(\varepsilon\); the mean VMAF change stays negligible and coverage is invariant to \(\varepsilon\).}
  \label{tab:abl:eps}
  \small
  \evaltable{% Auto-generated by ablation_eps_alpha.py
\begin{tabular}{lcccc}
\toprule
$\varepsilon$ (pts) & Bitrate saved (\%) & Rebuffer $\Delta$ (\%) & VMAF $\Delta$ (pts) & Coverage \\
\midrule
0.5 & 8.5 & -15.0 & +0.01 & 0.919 \\
1.0 & 16.5 & -35.5 & -0.01 & 0.919 \\
2.0 & 27.7 & -52.6 & -0.25 & 0.919 \\
4.0 & 38.4 & -67.1 & -0.42 & 0.919 \\
\bottomrule
\end{tabular}
}
\end{cps@table}

\begin{cps@table}
  \centering
  \caption{Sensitivity to the conformal miscoverage \(\alpha\) (greedy host, \mbox{5G} test, \(\varepsilon{=}\CPSepsilon\)). Empirical coverage tracks the finite-sample target \(1{-}\alpha\), and the intervention rate rises as \(\alpha\) tightens.}
  \label{tab:abl:alpha}
  \small
  \evaltable{% Auto-generated by ablation_eps_alpha.py
\begin{tabular}{lcccc}
\toprule
$\alpha$ & Target $1{-}\alpha$ & Coverage & Interv. (\%) & Bitrate saved (\%) \\
\midrule
0.05 & 0.95 & 0.945 & 73.0 & 4.5 \\
0.10 & 0.90 & 0.919 & 71.0 & 4.4 \\
0.20 & 0.80 & 0.854 & 35.2 & 6.2 \\
\bottomrule
\end{tabular}
}
\end{cps@table}

\begin{figure}[H]
  \centering
  \evalfig{figures/fig_cps_ablation.pdf}
  \caption{Hyperparameter sensitivity. Left: banked bitrate rises and the mean VMAF change stays near zero as the budget \(\varepsilon\) grows (\(\alpha{=}\CPSalpha\)). Right: empirical conformal coverage tracks the target \(1{-}\alpha\) across \(\alpha\) (\(\varepsilon{=}\CPSepsilon\)).}
  \label{fig:cps:ablation}
\end{figure}

% =========================================
\section{Discussion and limitations}
\label{sec:limits}

\textbf{Scope.}
CPS is an auditable, model-agnostic mechanism that (i)~banks bandwidth on VMAF-saturated ladder knees inside a stated perceptual budget and (ii)~enforces download feasibility under a conformal throughput lower bound with target coverage \(1{-}\alpha\).
These omissions were deliberate. The method did not claim dominance over every ABR controller, nor a pathwise worst-case guarantee beyond the conformal assumption of residual exchangeability.
The payoff also depended on the regime. Banking gains accrued where high rungs were often feasible (moderate-buffer synthetic \mbox{5G}). Under harsh broadband stress, the safety projection dominated and banking headroom decreased.

\textbf{Local ladder inversions.}
\tabref{tab:ladder:spacing} showed that, chunk by chunk, a sizable share of adjacent rungs in titles like Sintel and Big Buck Bunny ran near-flat or inverted. A higher bitrate then scored no better, and sometimes worse, in VMAF.
CPS handled all of this without a special case. Structurally, the evaluated shield ran on the monotone session-mean ladder, so Proposition~\ref{prop:collapse} applied outright. Even on a fully per-chunk, non-monotone ladder, both stages remained well defined. Each stage, moreover, handled inversions on its own terms. The feasibility projection~\eqref{eq:asafe} used download time, never VMAF, so inversions left safety untouched.
The knee rule~\eqref{eq:knee} picked the smallest rung within \(\varepsilon\) of the proposal's VMAF. When an inversion let a lower rung score at least as high, that rung satisfied the budget on its own, and banking captured the extra headroom rather than failing on it.
Local inversions, then, could only widen the banking opportunity, and the algorithm handled them by construction. The per-chunk ablation in Section~\ref{sec:eval:perchunk} confirmed this: banked bytes matched, and realized per-chunk perceptual cost stayed inside budget. Per-chunk ladders bought accuracy, not correctness.

\textbf{Threats to validity.}
\begin{enumerate}
  \item The broadband suite replayed real-measured FCC and Norway/Nornet traces~\cite{fcc2015mba,kvalbein2014}, yet the \mbox{5G} suite was synthetic. In both cases the chunk-level emulator omitted transport-stack dynamics and CDN caching~\cite{narayanan2021variegated}. Live-player validation was left to future work.
  \item Twelve publicly available titles and finite trace pools still capped catalogue diversity~\cite{kvalbein2014}; three were held out from PPO training but recurred in CPS evaluation.
  \item QoE here was a VMAF-based surrogate. Subjective rankings of stall against fidelity may diverge from it~\cite{seufert2015survey}.
  \item Conformal coverage presumed exchangeable residuals (Section~\ref{sec:method:conformal}). At \(n{=}204\) episodes, empirical coverage could sit marginally below the nominal \(1{-}\alpha\) target on real broadband traces (\(\CPSbbCov\) at \(\alpha{=}\CPSalpha\)) or at tight miscoverage on synthetic \mbox{5G} (\(0.945\) vs.\ \(0.95\) at \(\alpha{=}0.05\); \tabref{tab:abl:alpha}). Headline RQ2 reported the default \mbox{5G} operating point (\(\CPSGreedyCov\) vs.\ \(\CPScovTarget\)).
  \item Every arm---greedy, BBA, the PPO hosts---ran locally under the same finite-sample \((n{+}1)\) conformal correction. Co-design magnitudes hinged on the checkpoint, so we reported the best-validation model of each policy. The \(\Delta\) rows of \tabref{tab:cps:codesign} came from certified-arm means. The paired Wilcoxon test and bootstrap intervals used per-episode exports over the same \(\CPSepisodes\) pairs, with both co-design policies re-evaluated together. One point deserves emphasis: the co-design contrast used a single training seed per policy, so the \(95\%\) intervals spoke to evaluation variance, not training-seed variance. Establishing training-seed variance for \(w^{\star}\) would require retraining many seeds, which we left to future work. Hence co-design stood as a supporting result, not a headline claim.
  \item Banking savings scaled with ladder saturation (\tabref{tab:ladder:spacing}; see also the top-gain split in Section~\ref{sec:eval:main}). The twelve-title mix deliberately included steep titles so pooled effect sizes were conservative relative to a saturation-rich subset alone.
  \item The Pensieve-style no-op revealed that host aggressiveness set the banking opportunity, so results should not be carried over to arbitrary policies without the paired-arm protocol.
  \item The TOST margin equaled the configured banking budget \(\varepsilon\) (Section~\ref{sec:eval:setup}); we justified \(\varepsilon{=}\CPSepsilon\) relative to ladder spacing and perceptual encoding practice~\cite{menon2022perceptual}, but it remained a design parameter rather than a clinical equivalence margin~\cite{lakens2017equivalence}.
\end{enumerate}

\textbf{Future work.}
Several directions remain open: packet-level or field \mbox{5G} evaluation; applying CPS to RobustMPC and Comyco hosts; CVaR/tail-aware objectives~\cite{rockafellar2000}; telemetry-driven adaptation of \(\varepsilon\) and \(\alpha\); and preregistered non-inferiority studies with a fixed perceptual margin.
Real-player deployment was the highest-priority direction. Every result here came from a chunk-level emulator that omitted transport-stack dynamics. Inserting CPS into a production player such as \texttt{dash.js}~\cite{spiteri2018dashjs} therefore raised concrete challenges that the simulation never saw. Three challenges were central.
First, throughput had to be estimated from noisy, bursty HTTP/TCP (or QUIC) samples. Their exchangeability---the assumption behind the conformal bound---could break under slow-start, congestion-control transients, and CDN re-routing. The calibration residuals would then require periodic re-estimation, or a sliding, weighted-conformal variant that provably bounded coverage loss under drift~\cite{barber2023beyond}.
Second, the feasibility check had to reconcile the player's own buffer-occupancy signal, which is segment-level and delayed, with the certificate's buffer-margin accounting.
Third, the per-decision budget tightened from our \(\sim\!100\,\mu\mathrm{s}\) micro\-benchmark into a soft real-time constraint running beside the player's ABR loop.
Showing that the finite-sample coverage guarantee withstood these real-world perturbations---and measuring whatever coverage slack they opened up---was the natural next step.

% =========================================
\section{Conclusion}
\label{sec:conclusion}

We presented the Certified Perceptual Shield (CPS) for client-side ABR. It comprised three components: VMAF-knee bandwidth banking, an online conformal throughput lower bound with finite-sample coverage, and certified buffer feasibility.
Across \(\CPSepisodes\) paired synthetic \mbox{5G} episodes over twelve titles with diverse ladders, the results were consistent. Banking against conformal safety-only projection reduced delivered bitrate by about 4.5--3.8\% and rebuffering by about 9.3--9.2\% across every deterministic host (greedy through RobustMPC). Mean VMAF stayed inside \(\pm\CPSepsilon\), and empirical coverage was \(\CPSGreedyCov\) against a \(\CPScovTarget\) target.
Predictive banking further improved the margin by about 14.6\% greedy bitrate. Under broadband stress, safety alone predominated once the ladder knee lay beyond the observable operating range.
As a supporting single-seed result, optional PPO co-training with CPS in the loop was reported in Section~\ref{sec:eval:codesign} (\(+\CPScoDVM\) VMAF at \(+\CPScoDReb\)\,s rebuffering; crossover \(w^{\star}{=}\CPScoWstar\)).
CPS integrated with existing ABR controllers as a lightweight wrapper---no retraining required---and extended naturally to learned policies.

\appendix
% Appendix: Lagrangian CMDP training (optional RL hosts only).
\section{Lagrangian CMDP training for RL hosts}
\label{app:cmdp}

An unconstrained MDP maximizes \(\mathbb{E}[\sum_k \gamma^k r_k]\) alone.
When the base policy is learned, we optionally model streaming as a constrained MDP (CMDP) that maximizes expected fidelity subject to budgets on trajectory-level auxiliary costs~\cite{altman1999constrained}.
Let \(c^{\mathrm{stall}}_k \equiv \tau^{\mathrm{stall}}_k\) and \(c^{\mathrm{sm}}_k \equiv \text{smooth}_k\).
The undiscounted episodic rebuffering ratio and mean smoothness cost are
\begin{equation}
    \mathrm{RB}(e)=\frac{1}{T^{\mathrm{play}}_{e}}\sum_{k\in\mathcal{K}_{e}} c^{\mathrm{stall}}_k,
    \qquad
    \mathrm{SM}(e)=\frac{1}{K_{e}}\sum_{k\in\mathcal{K}_{e}} c^{\mathrm{sm}}_k,
\end{equation}
with targets \(\delta_{\mathrm{rf}}{=}0.05\) and \(\delta_{\mathrm{sm}}{=}3.5\)~\cite{yin2015control}.
Per-step normalized costs \(\widetilde{\mathrm{RB}}_k\) and \(\widetilde{\mathrm{SM}}_k\) yield the Lagrangian-relaxed PPO reward
\begin{equation}
\label{eq:rtrain-app}
    r^{\mathrm{train}}_k
    = \frac{1}{100}\Big(
        q_k
        - \lambda_r \,\widetilde{\mathrm{RB}}_k
        - \lambda_s \,\widetilde{\mathrm{SM}}_k
        - w_b\,\phi(B_k)
        - w_\ell\,\psi^{\mathrm{Lya}}_k
    \Big),
\end{equation}
with dual variables \((\lambda_r,\lambda_s)\) updated by projected ascent on empirical constraint gaps~\cite{paternain2019dual}.
Training optionally subtracts buffer-deviation and Lyapunov drift penalties~\cite{neely2010stochastic}; these stabilize optimization but are not formal certificates.
Finite-sample PPO does not certify per-chunk feasibility~\cite{schulman2017proximal,achiam2017constrained}; CPS supplies that layer at deployment time.

Implementation uses Stable-Baselines~3 PPO~\cite{raffin2021sb3,schulman2017proximal}: two Tanh MLPs (width 256), GAE \((\gamma{=}0.99,\lambda_{\mathrm{GAE}}{=}0.95)\), dual step sizes \(\eta_r{=}0.012\), \(\eta_s{=}0.003\), and projected ranges \(\lambda_r\in[4.3,40]\), \(\lambda_s\in[0.3,2.5]\) after a 60-episode warm-up.
Co-design wraps the training environment with CPS; shield-at-eval-only policies train without the wrapper.


% --- Back matter (Computer Networks / Elsevier: acknowledgements and declarations immediately before the reference list) ---
\section*{Data availability}
The code, evaluation scripts, and exported CSV files underlying \tabref{tab:cps:full}--\tabref{tab:cps:codesign} and \figref{fig:cps:overview}--\figref{fig:cps:codesign} are available with the article artifacts.
The Certified Perceptual Shield implementation and its parameters are documented with the code.

\section*{Funding}
This research did not receive any specific grant from funding agencies in the public, commercial, or not-for-profit sectors.

\section*{Declaration of competing interest}
The authors declare that they have no known competing financial interests or personal relationships that could have appeared to influence the work reported in this paper.

\section*{CRediT authorship contribution statement}
\textbf{Saeed~Zarbi:} Conceptualization, Methodology, Software, Validation, Formal analysis, Investigation, Data curation, Writing~-- original draft, Visualization.\par
\textbf{Leili~Farzinvash:} Supervision, Writing~-- review \& editing.\par
\textbf{Pedram~Salehpour:} Conceptualization, Supervision, Project administration, Writing~-- review \& editing.

% Flush deferred floats so the reference list is not interleaved with figures/tables.
\clearpage

% Elsevier numeric style (elsarticle bundle). If \texttt{elsarticle-num.bst} is missing, install the \texttt{elsarticle} CTAN package or use \texttt{unsrtnat} temporarily.
\bibliographystyle{elsarticle-num}
\bibliography{references}

\end{document}
